
% BEGIN S2
 
\par\medskip\noindent\textbf{2 Basic definitions and general results}\par
 
% BEGIN S2.p1
 
Given a signature $\Sigma$ consisting of a finite set of function symbols (with associated arities) and a countably infinite set of variables $V$, the set $T(\Sigma,V)$ of terms over $\Sigma$ with variables in $V$ is defined in the usual way [16]. An equational theory $E$ is given by a finite set of identities $s\approx t$ between terms, which are (implicitly) assumed to be universally quantified. Such a set of identities $E$ induces the congruence relation $\approx_{E}$ on terms, which can either be defined syntactically through rewriting or semantically through first-order interpretations of $\Sigma$, with $\approx$ as identity relation [16].\par
 
% END S2.p1
 
% BEGIN S2.p2
 
A substitution $\sigma$ is a mapping from $V$ to $T(\Sigma,V)$ that has a finite domain $\operatorname{Dom}(\sigma):=\{x\in V\mid\sigma(x)\neq x\}$. It can be homomorphically extended to a mapping from $T(\Sigma,V)$ to $T(\Sigma,V)$ by defining $\sigma(f(t_{1},\ldots,t_{n})):=f(\sigma(t_{1}),\ldots,\sigma(t_{n}))$. The variable range $\operatorname{VRan}(\sigma)$ of $\sigma$ consists of the set of variables occurring in the terms $\sigma(x)$ for $x\in\operatorname{Dom}(\sigma)$. Substitutions can be compared using the instantiation preorder: given an equational theory $E$, a set of variables $X\subseteq V$, and two substitutions $\sigma,\tau$, we say that $\sigma$ is more general than $\tau$ (or $\tau$ is an instance of $\sigma$) w.r.t. $E$ and $X$ (written $\sigma\leq_{E}^{X}\tau$) if there is a substitution $\lambda$ such that $\lambda\sigma\approx_{E}^{X}\tau$, i.e., $\lambda(\sigma(x))\approx_{E}\tau(x)$ holds for all $x\in X$. In case $X=V$ we also write $\lambda\sigma\approx_{E}\tau$ in place of $\lambda\sigma\approx_{E}^{V}\tau$. It is easy to see that $\leq_{E}^{X}$ is indeed a preorder, i.e., reflexive and transitive, but in general not antisymmetric. We write $\sim_{E}^{X}$ for the equivalence relation induced by $\leq_{E}^{X}$, i.e., $\sigma\sim_{E}^{X}\tau$ iff $\sigma\leq^{X}_{E}\tau$ and $\tau\leq^{X}_{E}\sigma$. We say that $\sigma$ is strictly more general than $\tau$ (or $\tau$ is a strict instance of $\sigma$) w.r.t. $E$ and $X$ (written $\sigma<_{E}^{X}\tau$) if $\sigma\leq_{E}^{X}\tau$ and $\sigma\not\sim_{E}^{X}\tau$.\par
 
% END S2.p2
 
% BEGIN S2.p3
 
An $E$-unification problem is a finite set of equations of the form $\Gamma=\{s_{1}\approx_{E}^{?}t_{1},\ldots,s_{n}\approx_{E}^{?}t_{n}\}$ such that $s_{1},t_{1},\ldots,s_{n},t_{n}$ are terms in $T(\Sigma,V)$. An $E$-unifier of $\Gamma$ is a substitution $\sigma$ that solves all the equations in $\Gamma$, i.e., satisfies $\sigma(s_{i})\approx_{E}\sigma(t_{i})$ for all $i,1\leq i\leq n$. The unification problem $\Gamma$ is solvable if it has an $E$-unifier. The set of all $E$-unifiers of $\Gamma$ is denoted as $\mathcal{U}_{E}(\Gamma)$. For elementary $E$-unification it is assumed that $\Sigma$ (and thus also $\Gamma$) contains only function symbols occurring in $E$. For $E$-unification with constants $\Sigma$ and $\Gamma$ may contain additional constant symbols, and for general $E$-unification $\Sigma$ and $\Gamma$ may contain additional function symbols of arbitrary arity.\par
 
% END S2.p3
 
% BEGIN S2.p4
 
Unification types for non-empty sets of identities $E$ are usually defined w.r.t. the restricted instantiation preorder, which is $\leq_{E}^{X}$ where $X$ is the finite set $\operatorname{Var}(\Gamma)$ of all variables occurring in the given unification problem $\Gamma$, but some authors also use the unrestricted instantiation preorder $\leq_{E}^{V}$. In this section, we will additionally consider settings where $X$ is between these two extremes. Note that $\operatorname{Var}(\Gamma)\subseteq X$ is required for the set of $E$-unifiers to be closed under instantiation.\par
 
% END S2.p4
 
% BEGIN Thmtheorem1
 
\par\medskip\noindent\textbf{Lemma 1.}\par
 
% BEGIN Thmtheorem1.p1
 
If $\Gamma$ is an $E$-unification problem and $X\subseteq V$ a set of variables satisfying $\operatorname{Var}(\Gamma)\subseteq X$, then $\sigma\in\mathcal{U}_{E}(\Gamma)$ implies $\theta\in\mathcal{U}_{E}(\Gamma)$ for all substitutions $\theta$ such that $\sigma\leq_{E}^{X}\theta$.\par
 
% END Thmtheorem1.p1
 
% END Thmtheorem1
 
% BEGIN S2.p5
 
Given an $E$-unification problem $\Gamma$ and some set of variables $X$ with $\operatorname{Var}(\Gamma)\subseteq X\subseteq V$, we say that a set $\mathcal{S}$ of substitutions is a complete set of $E$-unifiers of $\Gamma$ w.r.t. $\leq_{E}^{X}$ if it consists of $E$-unifiers of $\Gamma$, and every $E$-unifier of $\Gamma$ is an instance of an element of the complete set, i.e., for every $\theta\in\mathcal{U}_{E}(\Gamma)$ there exists $\sigma\in\mathcal{S}$ such that $\sigma\leq_{E}^{X}\theta$. Such a set is called minimal if it does not contain two distinct elements that are comparable w.r.t. $\leq_{E}^{X}$. It is easy to see that minimal complete sets of $E$-unifiers of a given unification problem $\Gamma$ are unique up to the equivalence relation $\sim_{E}^{X}$ induced by the preorder $\leq_{E}^{X}$ (see, e.g., Corollary 3.13 in [19] and Theorem 3 below), and thus all have the same cardinality.\par
 
% END S2.p5
 
% BEGIN Thmtheorem2
 
\par\medskip\noindent\textbf{Definition 2.}\par
 
% BEGIN Thmtheorem2.p1
 
Let $\Gamma$ be a solvable $E$-unification problem and $X$ a set of variables such that $\operatorname{Var}(\Gamma)\subseteq X\subseteq V$. Then the unification type of $\Gamma$ w.r.t. $\leq_{E}^{X}$ is\par
  
\par\noindent  $\blacksquare$ 
% BEGIN S2.I1.i1.p1
 
unitary if $\Gamma$ has a minimal complete set of $E$-unifiers of cardinality one w.r.t. $\leq_{E}^{X}$, whose single element is then called most general $E$-unifier (mgu),\par
 
% END S2.I1.i1.p1
 
 
\par\noindent  $\blacksquare$ 
% BEGIN S2.I1.i2.p1
 
finitary if $\Gamma$ has a finite minimal complete set of $E$-unifiers of cardinality greater than one w.r.t. $\leq_{E}^{X}$,\par
 
% END S2.I1.i2.p1
 
 
\par\noindent  $\blacksquare$ 
% BEGIN S2.I1.i3.p1
 
infinitary if $\Gamma$ has an infinite minimal complete set of $E$-unifiers w.r.t. $\leq_{E}^{X}$,\par
 
% END S2.I1.i3.p1
 
 
\par\noindent  $\blacksquare$ 
% BEGIN S2.I1.i4.p1
 
zero if $\Gamma$ does not have a minimal complete set of $E$-unifiers w.r.t. $\leq_{E}^{X}$, i.e., every complete set is redundant in the sense that it must contain two distinct elements that are comparable w.r.t. $\leq_{E}^{X}$.\par
 
% END S2.I1.i4.p1
 
  
% END Thmtheorem2.p1
 
% END Thmtheorem2
 
% BEGIN S2.p6
 
Minimal complete sets of unifiers can alternatively be characterized using the following order-theoretic point of view [3, 19]. Let $\Gamma$ be an $E$-unification problem and $X\subseteq V$ a set of variables satisfying $\operatorname{Var}(\Gamma)\subseteq X$. We denote the $\sim_{E}^{X}$-equivalence class of a unifier $\sigma$ as $[\sigma]_{E}^{X}$ and the set of all equivalence classes of unifiers as $[\mathcal{U}_{E}(\Gamma)]_{E}^{X}$. The partial order $\preceq_{E}^{X}$ on $[\mathcal{U}_{E}(\Gamma)]_{E}^{X}$ induced by the instantiation preorder $\leq_{E}^{X}$ on unifiers is defined as $[\sigma]_{E}^{X}\preceq_{E}^{X}[\tau]_{E}^{X}$ if $\sigma\leq_{E}^{X}\tau$. We say that $S\subseteq[\mathcal{U}_{E}(\Gamma)]_{E}^{X}$ is complete w.r.t. $\preceq_{E}^{X}$ if every element of $[\mathcal{U}_{E}(\Gamma)]_{E}^{X}$ is above (w.r.t. $\preceq_{E}^{X}$) some element of $S$.\par
 
% END S2.p6
 
% BEGIN Thmtheorem3
 
\par\medskip\noindent\textbf{Theorem 3 ([19]).}\par
 
% BEGIN Thmtheorem3.p1
 
Let $M$ be the set of $\preceq_{E}^{X}$-minimal elements of $[\mathcal{U}_{E}(\Gamma)]_{E}^{X}$. If $\mathcal{S}$ is a minimal complete set of $E$-unifiers of $\Gamma$ w.r.t. $\leq_{E}^{X}$, then $M=\{[\sigma]_{E}^{X}\mid\sigma\in\mathcal{S}\}$. Conversely, if $M$ is complete in $[\mathcal{U}_{E}(\Gamma)]_{E}^{X}$, then any set of substitutions obtained by picking one representative for each element of $M$ is a minimal complete set of $E$-unifiers of $\Gamma$.\par
 
% END Thmtheorem3.p1
 
% END Thmtheorem3
 
% BEGIN S2.p7
 
Consequently, unification type zero corresponds to the case where the set $M$ of minimal elements is not complete, whereas the other types are determined by the cardinality of the set $M$ in case it is complete. Theorem 3.1 in [3] establishes conditions that are necessary, sufficient, or both for proving unification type zero.333Note that, in [3], the instantiation preorder is written the other way round, i.e., more general substitutions are larger.  Here, we present one of these sufficient conditions since we will use it in Section 3 to show that the unification type of the description logic $\mathcal{E\mkern-1.618muL}$ is zero w.r.t. the restricted instantiation preorder. It was also employed in the first paper showing unification type zero for an equational theory [29, 30].\par
 
% END S2.p7
 
% BEGIN Thmtheorem4
 
\par\medskip\noindent\textbf{Lemma 4 ([3]).}\par
 
% BEGIN Thmtheorem4.p1
 
If there exists a strictly decreasing chain $\sigma_{1}>_{E}^{X}\sigma_{2}>_{E}^{X}\sigma_{3}>_{E}^{X}\cdots$ such that the set $\mathcal{S}=\{\sigma_{1},\sigma_{2},\sigma_{3},\ldots\}$ is a complete set of $E$-unifiers of $\Gamma$ w.r.t. $\leq_{E}^{X}$, then $\Gamma$ has type zero w.r.t. $\leq_{E}^{X}$.\par
 
% END Thmtheorem4.p1
 
% END Thmtheorem4
 
% BEGIN S2.SS0.SSS0.P0.SPx1
 
\par\medskip\noindent\textbf{Proof.}\par
 
% BEGIN S2.SS0.SSS0.P0.SPx1.p1
 
Suppose there exists a chain $\sigma_{1}>_{E}^{X}\sigma_{2}>_{E}^{X}\sigma_{3}>_{E}^{X}\cdots$ satisfying the conditions stated in the lemma. This chain satisfies the following properties:\par
  
\par\noindent  1. 
% BEGIN S2.I2.i1.p1
 
It has no lower bound in $\mathcal{U}_{E}(\Gamma)$, i.e., there is no $E$-unifier $\tau$ of $\Gamma$ such that $\sigma_{i}\geq_{E}^{X}\tau$ for all $i\geq 1$. To see why this is true, suppose such a unifier $\tau$ does exist. Since $\mathcal{S}$ is complete, there is $j\geq 1$ such that $\tau\geq_{E}^{X}\sigma_{j}>_{E}^{X}\sigma_{j+1}\geq_{E}^{X}\tau$. Transitivity of $\geq_{E}^{X}$ yields that $\sigma_{j+1}\geq_{E}^{X}\sigma_{j}$, but this contradicts $\sigma_{j}>_{E}^{X}\sigma_{j+1}$.\par
 
% END S2.I2.i1.p1
 
 
\par\noindent  2. 
% BEGIN S2.I2.i2.p1
 
For all $i\geq 1$, if there is $\tau\in\mathcal{U}_{E}(\Gamma)$ such that $\sigma_{i}\geq_{E}^{X}\tau$, then there exists $\tau^{\prime}\in\mathcal{U}_{E}(\Gamma)$ such that $\tau\geq_{E}^{X}\tau^{\prime}$ and $\sigma_{i+1}\geq_{E}^{X}\tau^{\prime}$. To show this, assume that $\sigma_{i}\geq_{E}^{X}\tau$. The completeness of $\mathcal{S}$ yields $j\geq 1$ such that $\sigma_{i}\geq_{E}^{X}\tau\geq_{E}^{X}\sigma_{j}$. This implies that $\sigma_{i}\geq_{E}^{X}\sigma_{j}$ because $\geq_{E}^{X}$ is a transitive relation. Hence, since the chain is strictly decreasing, it follows that $i\leq j$, and thus $\sigma_{i+1}\geq_{E}^{X}\sigma_{j+1}$ and $\tau\geq_{E}^{X}\sigma_{j+1}$. Consequently, we can set $\tau^{\prime}:=\sigma_{j+1}$.\par
 
% END S2.I2.i2.p1
 
  
Now, let $M$ be the set of $\preceq_{E}^{X}$-minimal elements of $[\mathcal{U}_{E}(\Gamma)]_{E}^{X}$. To prove that $\Gamma$ has type zero, it suffices to show that $M$ is not complete. Assume to the contrary that $M$ is complete. Then there exists $[\tau]^{X}_{E}\in M$ such that $[\tau]^{X}_{E}\preceq_{E}^{X}[\sigma_{1}]^{X}_{E}$, and thus $\sigma_{1}\geq_{E}^{X}\tau$. Since our chain does not have a lower bound, there must be an $i\geq 1$ such that $\sigma_{i}\geq_{E}^{X}\tau$ and $\sigma_{i+1}\not\geq_{E}^{X}\tau$. Using the second property shown for our chain, there exists a unifier $\tau^{\prime}$ such that $\tau\geq_{E}^{X}\tau^{\prime}$ and $\sigma_{i+1}\geq_{E}^{X}\tau^{\prime}$. Minimality of $[\tau]^{X}_{E}$ implies that $\tau\sim_{E}^{X}\tau^{\prime}$, and thus $\sigma_{i+1}\geq_{E}^{X}\tau^{\prime}\geq_{E}^{X}\tau$, which yields a contradiction. Thus, we have shown that $M$ cannot be complete, which implies that $\Gamma$ has type zero. $\hfill\blacktriangleleft$\par
 
% END S2.SS0.SSS0.P0.SPx1.p1
 
% BEGIN S2.SS0.SSS0.P0.SPx1.p2
 
As usual, we order unification types w.r.t. how bad they are (larger is worse) by setting\par
 
zero $>$ infinitary $>$ finitary $>$ unitary.\par
 
The unification type w.r.t. $\leq_{E}^{X}$ of an equational theory $E$ is the worst type of any solvable $E$-unification problem w.r.t. $\leq_{E}^{X}$. The unrestricted unification type of $\Gamma$ ($E$) is the one w.r.t. $\leq_{E}^{V}$ and the restricted unification type of $\Gamma$ ($E$) is the one w.r.t. $\leq_{E}^{X}$ for $X=\operatorname{Var}(\Gamma)$.\par
 
% END S2.SS0.SSS0.P0.SPx1.p2
 
% BEGIN S2.SS0.SSS0.P0.SPx1.p3
 
The results on the unification type of equational theories in the literature are usually shown for the restricted case. As we will demonstrate in this paper, it may indeed make a considerable difference for the unification type which instantiation preorder is employed in its definition. This is however not the case in the following situation.\par
 
% END S2.SS0.SSS0.P0.SPx1.p3
 
% BEGIN Thmtheorem5
 
\par\medskip\noindent\textbf{Lemma 5.}\par
 
% BEGIN Thmtheorem5.p1
 
Let $E$ be an equational theory, $\Gamma$ an $E$-unification problem, $X_{0}:=\operatorname{Var}(\Gamma)$, and $X\subseteq V$ a set of variables such that $X_{0}\subseteq X$ and $V\setminus X$ is infinite. If $\Gamma$ has a minimal complete set of $E$-unifiers w.r.t. $\leq_{E}^{X}$, then it has a minimal complete set of $E$-unifiers w.r.t. $\leq_{E}^{X_{0}}$ of the same cardinality, and vice versa.\par
 
% END Thmtheorem5.p1
 
% END Thmtheorem5
 
% END S2.SS0.SSS0.P0.SPx1
 
% BEGIN S2.SS0.SSS0.P0.SPx2
 
\par\medskip\noindent\textbf{Proof.}\par
 
% BEGIN S2.SS0.SSS0.P0.SPx2.p1
 
Let $\mathcal{S}$ be a minimal complete set of $E$-unifiers of $\Gamma$ w.r.t. $\leq_{E}^{X}$. For every unifier $\theta$ in $\mathcal{S}$ we can rename the variables in $\operatorname{VRan}(\theta)$ such that they do not belong to $X$ by applying an appropriate permutation $\pi$ that maps the variables of $\operatorname{VRan}(\theta)$ bijectively to a set of variables $Y$ with $|Y|=|\operatorname{VRan}(\theta)|$ and $Y\cap(X\cup\operatorname{VRan}(\theta))=\emptyset$. Such a finite set $Y$ of variables exists since $\operatorname{VRan}(\theta)$ is finite and $V\setminus X$ is infinite. For a given bijection $p:\operatorname{VRan}(\theta)\rightarrow Y$, we can define $\pi$ as follows: $\pi(z):=p(z)$ for all $z\in\operatorname{VRan}(\theta)$, $\pi(y):=p^{-1}(y)$ for all $y\in Y$, and $\pi(x)=x$ for all other variables. Note that $\pi$ is a substitution since its domain is $\operatorname{VRan}(\theta)\cup Y$, which is finite. To show that $\pi$ really is a permutation, i.e., a bijective mapping from $V$ to $V$, it is sufficient to show that it is an injective mapping from $V$ to $V$ (see Lemma 2.6 in [28]). This is an immediate consequence of the facts that $p$ and $p^{-1}$ are bijections and $Y$ and $\operatorname{VRan}(\theta)$ are disjoint.\par
 
% END S2.SS0.SSS0.P0.SPx2.p1
 
% BEGIN S2.SS0.SSS0.P0.SPx2.p2
 
The substitution $\pi\theta$ is equivalent to $\theta$ w.r.t. the equivalence relation $\sim_{\emptyset}^{V}$ induced by $\leq_{\emptyset}^{V}$, and thus also w.r.t. $\sim_{E}^{X}$. In fact, $\pi^{-1}\pi\theta=\theta$. If we restrict the domain of this substitution to $X$, then the resulting substitution $\pi\theta|_{X}$ is still equivalent to $\theta$ w.r.t. $\sim_{E}^{X}$ and also satisfies $\operatorname{VRan}(\pi\theta|_{X})\cap X=\emptyset$. To see the latter, consider a variable $x\in\operatorname{Dom}(\pi\theta|_{X})$. First assume that $x\in\operatorname{Dom}(\theta)\cap X$. Then all the variables in $\theta(x)$ belong to $\operatorname{VRan}(\theta)$, and thus all the variables in $\pi\theta|_{X}(x)$ belong to $Y$, which is disjoint with $X$. If $x\in X\setminus\operatorname{Dom}(\theta)$, then $\pi\theta|_{X}(x)=\pi(x)\neq x$. Since $\operatorname{Dom}(\pi)=\operatorname{VRan}(\theta)\cup Y$ and $Y$ is disjoint with $X$, this implies $x\in\operatorname{VRan}(\theta)$, and thus $\pi(x)\in Y$.\par
 
% END S2.SS0.SSS0.P0.SPx2.p2
 
% BEGIN S2.SS0.SSS0.P0.SPx2.p3
 
Let $\mathcal{S}^{\prime}$ be the set of substitutions obtained from $\mathcal{S}$ by applying this renaming and domain restriction process to every element of $\mathcal{S}$. Due to the $\sim_{E}^{X}$-equivalence of the elements of $\mathcal{S}$ with their modified variants in $\mathcal{S}^{\prime}$ and the fact that $X$ contains all the variables occurring in $\Gamma$, it is easy to see that $\mathcal{S}^{\prime}$ is also a minimal complete set of $E$-unifiers of $\Gamma$ w.r.t. $\leq_{E}^{X}$. In a second step, we restrict the domains of the elements of $\mathcal{S}^{\prime}$ to $X_{0}$. Given an element $\theta$ of $\mathcal{S}^{\prime}$, we define $\theta|_{X_{0}}$ to be the substitution that coincides with $\theta$ on $X_{0}$ and maps each variable $x\not\in X_{0}$ to $x$. Since $\theta$ is an $E$-unifier of $\Gamma$ and this unification problem contains only variables from $X_{0}$, the substitution $\theta|_{X_{0}}$ is clearly also an $E$-unifier of $\Gamma$. We claim that $\theta|_{X_{0}}\leq_{E}^{X}\theta$. To prove this, we define the substitution $\lambda$ by setting $\lambda(x):=\theta(x)$ if $x\in X\setminus X_{0}$ and $\lambda(x):=x$ for all other variables. We now show that $\lambda\theta|_{X_{0}}\approx_{E}^{X}\theta$. If $x\in X_{0}\cap\operatorname{Dom}(\theta)$, then $\lambda(\theta|_{X_{0}}(x))=\lambda(\theta(x))=\theta(x)$. The first identity holds by the definition of $\theta|_{X_{0}}$ and the second since the variables occurring in $\theta(x)$ are elements of $\operatorname{VRan}(\theta)$, and thus do not belong to $X$. If $x\in X_{0}\setminus\operatorname{Dom}(\theta)$, then $\lambda(\theta|_{X_{0}}(x))=\lambda(x)=x=\theta(x)$. Finally, if $x\in X\setminus X_{0}$, then $\lambda(\theta|_{X_{0}}(x))=\lambda(x)=\theta(x)$. Again, the first identity holds by the definition of $\theta|_{X_{0}}$ and the second by the definition of $\lambda$. Summing up, we have shown that the following holds for every element $\theta$ of $\mathcal{S}^{\prime}$: $\theta|_{X_{0}}$ is an $E$-unifier of $\Gamma$ and $\theta|_{X_{0}}\leq_{E}^{X}\theta$. Since $\mathcal{S}^{\prime}$ is a minimal complete set of $E$-unifiers of $\Gamma$ w.r.t. $\leq_{E}^{X}$, its elements are minimal w.r.t. $\leq_{E}^{X}$, which yields $\theta|_{X_{0}}\sim_{E}^{X}\theta$. Consequently, we know that the set $\mathcal{S}^{\prime}|_{X_{0}}:=\{\theta|_{X_{0}}\mid\theta\in\mathcal{S}^{%
\prime}\}$ is also a minimal complete set of $E$-unifiers of $\Gamma$ w.r.t. $\leq_{E}^{X}$.\par
 
% END S2.SS0.SSS0.P0.SPx2.p3
 
% BEGIN S2.SS0.SSS0.P0.SPx2.p4
 
We claim that the same is true w.r.t. the smaller set of variables $X_{0}$, i.e., that $\mathcal{S}^{\prime}|_{X_{0}}$ is also minimal and complete w.r.t. $\leq_{E}^{X_{0}}$. Completeness trivially follows from the fact that ${\leq_{E}^{X}}\subseteq{\leq_{E}^{X_{0}}}$. To prove minimality, assume that $\theta$ and $\tau$ are two distinct elements of $\mathcal{S}^{\prime}|_{X_{0}}$. Then these two substitutions are not comparable w.r.t. $\leq_{E}^{X}$. Assume that they are comparable w.r.t. $\leq_{E}^{X_{0}}$, i.e., there is a substitution $\lambda$ such that $\lambda(\theta(x))\approx_{E}\tau(x)$ holds for all $x\in X_{0}$. By our construction of $\mathcal{S}^{\prime}|_{X_{0}}$, we know that $\operatorname{Dom}(\theta)\subseteq X_{0}$, $\operatorname{Dom}(\tau)\subseteq X_{0}$, and $\operatorname{VRan}(\theta)\cap X=\emptyset$. If $x$ is a variable in $X\setminus X_{0}$, then $\lambda(\theta(x))=\lambda(x)$ and $\tau(x)=x$. Thus, if we modify $\lambda$ to $\lambda^{\prime}$ such that $\lambda^{\prime}(x)=x$ holds for all $x\in X\setminus X_{0}$, then $\lambda^{\prime}(\theta(x))=\tau(x)$ holds for all $x\in X\setminus X_{0}$. This modification has no effect on the variables $x\in X_{0}$. In fact, let $x$ be such a variable. If $x\in\operatorname{Dom}(\theta)$, then $\theta(x)$ does not contain any variable from $X$, and thus $\lambda^{\prime}(\theta(x))=\lambda(\theta(x))\approx_{E}\tau(x)$ holds. If $x\notin\operatorname{Dom}(\theta)$, then $\theta(x)=x$, and thus again $\lambda^{\prime}(\theta(x))=\lambda^{\prime}(x)=\lambda(x)=\lambda(\theta(x))%
\approx_{E}\tau(x)$ since $\lambda^{\prime}$ coincides with $\lambda$ on the variables in $X_{0}$. Summing up, we have shown that the assumption $\theta\leq_{E}^{X_{0}}\tau$ implies $\theta\leq_{E}^{X}\tau$, which contradicts the fact that $\mathcal{S}^{\prime}|_{X_{0}}$ is minimal w.r.t. $\leq_{E}^{X}$. Consequently, $\mathcal{S}^{\prime}|_{X_{0}}$ is a minimal complete set w.r.t. $\leq_{E}^{X_{0}}$, and this set has the same cardinality as the minimal complete set $\mathcal{S}$ of $E$-unifiers of $\Gamma$ w.r.t. $\leq_{E}^{X}$ we have started with.\par
 
% END S2.SS0.SSS0.P0.SPx2.p4
 
% BEGIN S2.SS0.SSS0.P0.SPx2.p5
 
Conversely, let $\mathcal{S}$ be a minimal complete set of $E$-unifiers of $\Gamma$ w.r.t. $\leq_{E}^{X_{0}}$. By applying the same construction as in the proof of the other direction, we can assume without loss of generality that every unifier $\theta\in\mathcal{S}$ satisfies $\operatorname{Dom}(\theta)\subseteq X_{0}$ and $\operatorname{VRan}(\theta)\cap X=\emptyset$. In fact, by applying this construction to a minimal complete set $\mathcal{S}^{\prime}$ of $E$-unifiers of $\Gamma$ w.r.t. $\leq_{E}^{X_{0}}$, we obtain a new set $\mathcal{S}$ where every element $\theta^{\prime}$ of the original set $\mathcal{S}^{\prime}$ is replaced with an element $\theta$ that satisfies the above conditions and is $\sim_{E}^{X}$-equivalent to $\theta^{\prime}$, and thus also $\sim_{E}^{X_{0}}$-equivalent. Consequently, this new set $\mathcal{S}$ is also a minimal complete set of $E$-unifiers of $\Gamma$ w.r.t. $\leq_{E}^{X_{0}}$.\par
 
% END S2.SS0.SSS0.P0.SPx2.p5
 
% BEGIN S2.SS0.SSS0.P0.SPx2.p6
 
We claim that $\mathcal{S}$ is also minimal and complete w.r.t. the larger set of variables $X$. Minimality trivially follows from the fact that ${\leq_{E}^{X}}\subseteq{\leq_{E}^{X_{0}}}$. To prove completeness, assumed that $\tau$ is an $E$-unifier of $\Gamma$. Completeness of $\mathcal{S}$ w.r.t. $\leq_{E}^{X_{0}}$ yields an element $\theta$ of $\mathcal{S}$ such that $\theta\leq_{E}^{X_{0}}\tau$, i.e., there is a substitution $\lambda$ such that $\lambda\theta(x)\approx_{E}\tau(x)$ holds for all $x\in X_{0}$. Let $x$ be a variable in $X\setminus X_{0}$. Then $\lambda\theta(x)=\lambda(x)$. Thus, if we modify $\lambda$ to $\lambda^{\prime}$ such that $\lambda^{\prime}(x)=\tau(x)$ holds for all $x\in X\setminus X_{0}$, then $\lambda^{\prime}(\theta(x))=\tau(x)$ holds for all $x\in X\setminus X_{0}$. The claim that this modification has no effect for the variables $x\in X_{0}$ can be shown as in the proof of minimality of $\mathcal{S}^{\prime}|_{X_{0}}$ w.r.t. $\leq_{E}^{X_{0}}$ above. This proves that $\lambda^{\prime}\theta(x)\approx_{E}\tau(x)$ holds for all $x\in X$, and thus $\mathcal{S}$ is also complete w.r.t. $\leq_{E}^{X}$. $\hfill\blacktriangleleft$\par
 
% END S2.SS0.SSS0.P0.SPx2.p6
 
% BEGIN S2.SS0.SSS0.P0.SPx2.p7
 
The following theorem is an immediate consequence of this lemma.\par
 
% END S2.SS0.SSS0.P0.SPx2.p7
 
% BEGIN Thmtheorem6
 
\par\medskip\noindent\textbf{Theorem 6.}\par
 
% BEGIN Thmtheorem6.p1
 
Let $E$ be an equational theory, $\Gamma$ an $E$-unification problem, and $X\subseteq V$ a set of variables such that $\operatorname{Var}(\Gamma)\subseteq X$ and $V\setminus X$ is infinite. Then the restricted unification type of $\Gamma$ coincides with the unification type of $\Gamma$ w.r.t. $\leq_{E}^{X}$.\par
 
% END Thmtheorem6.p1
 
% END Thmtheorem6
 
% BEGIN S2.SS0.SSS0.P0.SPx2.p8
 
Note that the condition of the theorem is in particular satisfied if $X$ is a finite superset of $\operatorname{Var}(\Gamma)$. However, it is clearly not satisfied for $X=V$, which corresponds to the unrestricted instantiation preorder setting. We will see below that in this case the unification type can indeed depend on whether the restricted or the unrestricted instantiation preorder is used. However, there is a special case where this cannot happen.\par
 
% END S2.SS0.SSS0.P0.SPx2.p8
 
% BEGIN Thmtheorem7
 
\par\medskip\noindent\textbf{Theorem 7.}\par
 
% BEGIN Thmtheorem7.p1
 
Let $E$ be an equational theory, $\Gamma$ an $E$-unification problem, and $\mathcal{S}$ a set of $E$-unifiers of $\Gamma$ such that $\operatorname{VRan}(\sigma)\cup\operatorname{Dom}(\sigma)\subseteq%
\operatorname{Var}(\Gamma)$ holds for all $\sigma\in\mathcal{S}$. Then $\mathcal{S}$ is a minimal complete set of $E$-unifiers of $\Gamma$ w.r.t. $\leq_{E}^{V}$ iff it is a minimal complete set of $E$-unifiers of $\Gamma$ w.r.t. $\leq_{E}^{\operatorname{Var}(\Gamma)}$.\par
 
% END Thmtheorem7.p1
 
% END Thmtheorem7
 
% END S2.SS0.SSS0.P0.SPx2
 
% BEGIN S2.SS0.SSS0.P0.SPx3
 
\par\medskip\noindent\textbf{Proof.}\par
 
% BEGIN S2.SS0.SSS0.P0.SPx3.p1
 
Let $\mathcal{S}$ be a minimal complete set of $E$-unifiers of $\Gamma$ w.r.t. $\leq_{E}^{V}$. Since ${\leq_{E}^{V}}\subseteq{\leq_{E}^{\operatorname{Var}(\Gamma)}}$, this set is also complete w.r.t. $\leq_{E}^{\operatorname{Var}(\Gamma)}$. To show minimality w.r.t. $\leq_{E}^{\operatorname{Var}(\Gamma)}$, assume to the contrary that $\sigma,\theta$ are two distinct elements of $\mathcal{S}$ such that $\sigma\leq_{E}^{\operatorname{Var}(\Gamma)}\theta$, i.e., there is a substitution $\lambda$ such that $\lambda\sigma(x)\approx_{E}\theta(x)$ holds for all $x\in\operatorname{Var}(\Gamma)$. We modify $\lambda$ to $\lambda^{\prime}$ by setting $\lambda^{\prime}(x)=x$ for all variables $x\in V\setminus\operatorname{Var}(\Gamma)$. For $x\in\operatorname{Var}(\Gamma)$, we know that $\sigma(x)$ contains only variables from $\operatorname{VRan}(\sigma)\subseteq\operatorname{Var}(\Gamma)$ if $x\in\operatorname{Dom}(\sigma)$ or $\sigma(x)=x\in\operatorname{Var}(\Gamma)$. Since $\lambda$ and $\lambda^{\prime}$ coincide on $\operatorname{Var}(\Gamma)$, this yields $\lambda^{\prime}\sigma(x)=\lambda\sigma(x)\approx_{E}\theta(x)$. For $x\in V\setminus\operatorname{Var}(\Gamma)$, this variable does not belong to any of the sets $\operatorname{Dom}(\sigma)$, $\operatorname{Dom}(\theta)$, and $\operatorname{Dom}(\lambda^{\prime})$, and thus $\lambda^{\prime}\sigma(x)=\lambda^{\prime}(x)=x=\theta(x)$. Summing up, we have shown that $\sigma\leq_{E}^{V}\theta$, which contradicts our assumption that $\mathcal{S}$ is minimal w.r.t. $\leq_{E}^{V}$.\par
 
% END S2.SS0.SSS0.P0.SPx3.p1
 
% BEGIN S2.SS0.SSS0.P0.SPx3.p2
 
Conversely, assume that $\mathcal{S}$ is a minimal complete set of $E$-unifiers of $\Gamma$ w.r.t. $\leq_{E}^{\operatorname{Var}(\Gamma)}$. Since ${\leq_{E}^{V}}\subseteq{\leq_{E}^{\operatorname{Var}(\Gamma)}}$, minimality also holds w.r.t. $\leq_{E}^{V}$. To show completeness w.r.t. $\leq_{E}^{V}$, assume that $\theta$ is an $E$-unifier of $\Gamma$. Completeness of $\mathcal{S}$ w.r.t. $\leq_{E}^{\operatorname{Var}(\Gamma)}$ yields a substitution $\sigma\in\mathcal{S}$ such that $\sigma\leq_{E}^{\operatorname{Var}(\Gamma)}\theta$. Similarly to the first part of the proof, we can show that this also implies $\sigma\leq_{E}^{V}\theta$. The difference is that now $\theta$ is an arbitrary unifier, and thus $\operatorname{VRan}(\theta)\cup\operatorname{Dom}(\theta)\subseteq%
\operatorname{Var}(\Gamma)$ need not hold. Let $\lambda$ be such that $\lambda\sigma(x)\approx_{E}\theta(x)$ holds for all $x\in\operatorname{Var}(\Gamma)$. We modify $\lambda$ to $\lambda^{\prime}$ by setting $\lambda^{\prime}(x)=\theta(x)$ for all variables $x\in V\setminus\operatorname{Var}(\Gamma)$. For $x\in\operatorname{Var}(\Gamma)$, we obtain $\lambda^{\prime}\sigma(x)=\lambda\sigma(x)\approx_{E}\theta(x)$ as in the first part of the proof. For $x\in V\setminus\operatorname{Var}(\Gamma)$, this variable does not belong to $\operatorname{Dom}(\sigma)$, and thus $\lambda^{\prime}\sigma(x)=\lambda^{\prime}(x)=\theta(x)$. $\hfill\blacktriangleleft$\par
 
% END S2.SS0.SSS0.P0.SPx3.p2
 
% BEGIN S2.SS0.SSS0.P0.SPx3.p3
 
Examples of equational theories where the conditions of this theorem are always satisfied are the empty theory (unitary), the theory $\mathsf{C}$ axiomatizing commutativity of a binary function symbol (finitary), and the theory $\mathsf{A}$ axiomatizing associativity of a binary function symbol (infinitary). This is an easy consequence of the known algorithms [42, 50, 45] computing (or enumerating, in the case of $\mathsf{A}$) minimal complete sets of unifiers for these theories.\par
 
% END S2.SS0.SSS0.P0.SPx3.p3
 
% BEGIN S2.SS0.SSS0.P0.SPx3.p4
 
The following result is an easy consequence of the fact that ${\leq_{E}^{V}}\subseteq{\leq_{E}^{\operatorname{Var}(\Gamma)}}$.\par
 
% END S2.SS0.SSS0.P0.SPx3.p4
 
% BEGIN Thmtheorem8
 
\par\medskip\noindent\textbf{Theorem 8.}\par
 
% BEGIN Thmtheorem8.p1
 
Let $E$ be an equational theory. If $E$ has unrestricted unification type unitary (finitary), then it has restricted unification type unitary (finitary or unitary).\par
 
% END Thmtheorem8.p1
 
% END Thmtheorem8
 
% END S2.SS0.SSS0.P0.SPx3
 
% BEGIN S2.SS0.SSS0.P0.SPx4
 
\par\medskip\noindent\textbf{Proof.}\par
 
% BEGIN S2.SS0.SSS0.P0.SPx4.p1
 
A given finite minimal complete set $\mathcal{S}$ of $E$-unifiers of $\Gamma$ w.r.t. $\leq_{E}^{V}$ is also complete w.r.t. $\leq_{E}^{\operatorname{Var}(\Gamma)}$. If it is not minimal (which can only happen in the finitary case), then one can make it minimal by removing redundant element (i.e., elements $\theta$ such that $\mathcal{S}$ contains an element $\sigma<_{E}^{\operatorname{Var}(\Gamma)}\theta$) without destroying completeness. $\hfill\blacktriangleleft$\par
 
% END S2.SS0.SSS0.P0.SPx4.p1
 
% END S2.SS0.SSS0.P0.SPx4
 
% END S2


% BEGIN S3
 
\par\medskip\noindent\textbf{3 The unification type of $\mathcal{E\mkern-1.618muL}$}\par
 
% BEGIN S3.p1
 
Unification was introduced in description logics as a tool for detecting redundancies in large knowledge bases [14]. The description logic $\mathcal{E\mkern-1.618muL}$ [6] has drawn considerable attention since its standard reasoning problems can be solved in polynomial time while the logic is still expressive enough for formalizing bio-medical ontologies [11]. In [12, 13] it was shown that, in the setting of unification with constants, $\mathcal{E\mkern-1.618muL}$ has unification type zero w.r.t. the restricted instantiation preorder. In the following, we will analyze the example used in [13] to prove this result in more detail, and explain why it does not work for the unrestricted instantiation preorder. Then we show that the unification type of $\mathcal{E\mkern-1.618muL}$ in the unrestricted setting is actually infinitary, not just for this example, but in general.\par
 
% END S3.p1
 
% BEGIN S3.p2
 
But first, we briefly introduce $\mathcal{E\mkern-1.618muL}$ and recall why unification in $\mathcal{E\mkern-1.618muL}$ can be seen as unification modulo an equational theory (see [6, 13, 15] for more detailed descriptions). Given sets of concept names (unary predicates) and role names (binary predicates), $\mathcal{E\mkern-1.618muL}$ concept descriptions (or simply concepts) are built from concept names using the concept constructors top concept ($\top$), conjunction ($C\sqcap D$), and existential restriction ($\exists r.C$). In the model-theoretic semantics of $\mathcal{E\mkern-1.618muL}$, a given interpretation $\mathcal{I}$ assigns sets $C^{\mathcal{I}}$ to concept descriptions $C$ according to the semantics of the constructors. To be more precise, an interpretation $\mathcal{I}=(\Delta^{\mathcal{I}},\cdot^{\mathcal{I}})$ consists of a non-empty interpretation domain and an extension function that assigns subsets of this domain to concept names and binary relations on the domain to role names. This interpretation function is extended to concept descriptions as follows:\par
    \[\top^{\mathcal{I}}:=\Delta^{\mathcal{I}},\ \ (C\sqcap D)^{\mathcal{I}}:=C^{%
\mathcal{I}}\cap D^{\mathcal{I}},\ \ (\exists r.C)^{\mathcal{I}}:=\{d\in\Delta%
^{\mathcal{I}}\mid\exists e\in C^{\mathcal{I}}\ \mbox{such that}\ (d,e)\in r^{%
\mathcal{I}}\}.\]    
The concept description $C$ is subsumed by the concept description $D$ (written $C\sqsubseteq D$) if $C^{\mathcal{I}}\subseteq D^{\mathcal{I}}$ holds for all interpretations $\mathcal{I}$, and $C$ and $D$ are equivalent (written $C\equiv D$) if they subsume each other.\par
 
% END S3.p2
 
% BEGIN S3.p3
 
In [13], the notion of a reduced $\mathcal{E\mkern-1.618muL}$ concept is employed to derive characterizations of equivalence and subsumption of $\mathcal{E\mkern-1.618muL}$ concepts. Here, we recall this notion and the characterizations since they are used later on to explain why $\mathcal{E\mkern-1.618muL}$ has unification type zero w.r.t. the restricted instantiation preorder. Küsters [39] introduces the following rules for reducing $\mathcal{E\mkern-1.618muL}$ concept descriptions:\par
    \[\begin{array}[]{l@{\ \ \ \ }l}C\sqcap\top\rightarrow C&\mbox{for all $\mathcal%
{EL}$ concept descriptions $C$},\\[3.00003pt]
A\sqcap A\rightarrow A&\mbox{for all concept names $A$},\\[3.00003pt]
\exists r.C\sqcap\exists r.D\rightarrow\exists r.C&\mbox{for all $\mathcal{EL}%
$ concept descriptions $C,D$
with $C\sqsubseteq D$}.\end{array}\]    
A reduced form of a given $\mathcal{E\mkern-1.618muL}$ concept description $C$ is then obtained from $C$ by applying these rules exhaustively modulo associativity and commutativity of $\sqcap$. The following theorem is stated in [13] as an easy consequence of Corollary 6.3.1 on page 181 of [39].\par
 
% END S3.p3
 
% BEGIN Thmtheorem9
 
\par\medskip\noindent\textbf{Theorem 9 (Theorem 3.1 of [13]).}\par
 
% BEGIN Thmtheorem9.p1
 
Let $C,D$ be $\mathcal{EL}$ concept descriptions, and $\widehat{C},\widehat{D}$ reduced forms of $C,D$, respectively. Then $C\equiv D$ iff $\widehat{C}$ is identical to $\widehat{D}$ up to associativity and commutativity of $\sqcap$.\par
 
% END Thmtheorem9.p1
 
% END Thmtheorem9
 
% BEGIN S3.p4
 
This theorem is used in [13] to derive a recursive characterization of subsumption in $\mathcal{EL}$.\par
 
% END S3.p4
 
% BEGIN Thmtheorem10
 
\par\medskip\noindent\textbf{Corollary 10 (Corollary 3.2 of [13]).}\par
 
% BEGIN Thmtheorem10.p1
 
Let $C=A_{1}\sqcap\ldots\sqcap A_{k}\sqcap\exists r_{1}.C_{1}\sqcap\ldots\sqcap%
\exists r_{m}.C_{m}$ and $D=B_{1}\sqcap\ldots\sqcap B_{\ell}\sqcap\exists s_{1}.D_{1}\sqcap\ldots\sqcap%
\exists s_{n}.D_{n}$, where $A_{1},\ldots,A_{k},B_{1},\ldots,B_{\ell}$ are concept names. Then $C\sqsubseteq D$ iff $\{B_{1},\ldots,B_{\ell}\}\subseteq\{A_{1},\ldots,A_{k}\}$ and for every $j,1\leq j\leq n$, there exists an $i,1\leq i\leq m$, such that $r_{i}=s_{j}$ and $C_{i}\sqsubseteq D_{j}$.\par
 
% END Thmtheorem10.p1
 
% END Thmtheorem10
 
% BEGIN S3.p5
 
From this corollary, the following lemma is then derived in [13].\par
 
% END S3.p5
 
% BEGIN Thmtheorem11
 
\par\medskip\noindent\textbf{Lemma 11 (Lemma 3.3 of [13]).}\par
 
% BEGIN Thmtheorem11.p1
 
If $C,D$ are reduced $\mathcal{EL}$ concept descriptions such that $\exists r.D\sqsubseteq C$, then $C$ is either $\top$, or of the form $C=\exists r.C_{1}\sqcap\ldots\sqcap\exists r.C_{n}$ where $n\geq 1$; $C_{1},\ldots,C_{n}$ are reduced and pairwise incomparable w.r.t. subsumption; and $D\sqsubseteq C_{1},\ldots,D\sqsubseteq C_{n}$. Conversely, if $C,D$ are $\mathcal{EL}$ concept descriptions such that $C=\exists r.C_{1}\sqcap\ldots\sqcap\exists r.C_{n}$ and $D\sqsubseteq C_{1},\ldots,D\sqsubseteq C_{n}$, then $\exists r.D\sqsubseteq C$.\par
 
% END Thmtheorem11.p1
 
% END Thmtheorem11
 
% BEGIN S3.p6
 
Equivalence of $\mathcal{E\mkern-1.618muL}$ concept descriptions can be axiomatized by the equational theory $\mathsf{bSLmO}$ of bounded semilattices with monotone operators [54, 15, 55]. For this purpose, we view the conjunction operator $\sqcap$ as a binary function symbol, $\top$ as a constant symbol, and $\exists r.$ for a role name $r$ as a unary function symbol. The theory $\mathsf{bSLmO}$ over this signature then consists of the identities stating that $\sqcap$ is associative, commutative, and idempotent, has $\top$ as unit, and existential restrictions as monotone operators:\par
    \[x\sqcap y\approx y\sqcap x,\ (x\sqcap y)\sqcap z\approx x\sqcap(y\sqcap z),\ x%
\sqcap x\approx x,\ x\sqcap\top\approx x,\exists r.x\sqcap\exists r.(x\sqcap y%
)\approx\exists r.(x\sqcap y).\]    
% END S3.p6
 
% BEGIN S3.p7
 
For unification in $\mathcal{E\mkern-1.618muL}$ we partition the set of concept names into concept variables and concept constants. Substitutions can then replace concept variables in $\mathcal{E\mkern-1.618muL}$ concept descriptions with $\mathcal{E\mkern-1.618muL}$ concept descriptions. Unifiers are supposed to solve equations between concept descriptions with variables by making them equivalent. This corresponds to unification with constants modulo the equational theory $\mathsf{bSLmO}$. In the following, we use $\approx_{\mathcal{E\mkern-1.618muL}}$ (rather than $\approx_{\mathsf{bSLmO}}$) to denote equivalence between $\mathcal{E\mkern-1.618muL}$ concept descriptions and also employ the subscript $\mathcal{E\mkern-1.618muL}$ when writing the respective instantiation preorders.\par
 
% END S3.p7
 
% BEGIN Thmtheorem12
 
\par\medskip\noindent\textbf{Example 12.}\par
 
% BEGIN Thmtheorem12.p1
 
Consider the $\mathcal{E\mkern-1.618muL}$-unification problem $\Gamma:=\{x\sqcap\exists r.y\approx_{\mathcal{E\mkern-1.618muL}}^{?}\exists r.y\}$. Then, for every $n\geq 0$, the substitution\par
    \[\sigma_{n}:=\{x\mapsto\exists r.z_{1}\sqcap\ldots\sqcap\exists r.z_{n},\ \ y%
\mapsto z_{1}\sqcap\ldots\sqcap z_{n}\sqcap z\}\]    
for distinct variables $x,y,z_{1},\ldots,z_{n},z$ is an $\mathcal{E\mkern-1.618muL}$-unifier of $\Gamma$, where the empty conjunction is $\top$, i.e., $\sigma_{0}=\{x\mapsto\top,y\mapsto z\}$. We will show below that, w.r.t. the restricted instantiation preorder, the set $\{\sigma_{n}\mid n\geq 0\}$ is a complete set of $\mathcal{E\mkern-1.618muL}$-unifiers that constitutes a strictly decreasing chain $\sigma_{0}>_{\mathcal{E\mkern-1.618muL}}^{\{x,y\}}\sigma_{1}>_{\mathcal{E%
\mkern-1.618muL}}^{\{x,y\}}\sigma_{2}>_{\mathcal{E\mkern-1.618muL}}^{\{x,y\}}\cdots$ of more and more general unifiers. According to Lemma 4 this implies that $\Gamma$ has unification type zero, and thus $\mathcal{E\mkern-1.618muL}$ has unification type zero.\par
 
% END Thmtheorem12.p1
 
% END Thmtheorem12
 
% BEGIN Thmtheorem13
 
\par\medskip\noindent\textbf{Lemma 13.}\par
 
% BEGIN Thmtheorem13.p1
 
Let $\Gamma$ and $\sigma_{n}$ for $n\geq 0$ be defined as in Example 12. Then the set $\{\sigma_{n}\mid n\geq 0\}$ is a complete set of $\mathcal{E\mkern-1.618muL}$-unifiers of $\Gamma$ w.r.t. $\leq_{\mathcal{E\mkern-1.618muL}}^{\{x,y\}}$.\par
 
% END Thmtheorem13.p1
 
% END Thmtheorem13
 
% BEGIN S3.SS0.SSS0.P0.SPx1
 
\par\medskip\noindent\textbf{Proof.}\par
 
% BEGIN S3.SS0.SSS0.P0.SPx1.p1
 
Let $\sigma$ be an $\mathcal{E\mkern-1.618muL}$-unifier of $\Gamma$ and $C:=\sigma(y)$ and $D:=\sigma(x)$. Then $\exists r.C\sqsubseteq D$. In addition, we can assume without loss of generality that $C$ and $D$ are reduced. The characterization of subsumption in $\mathcal{E\mkern-1.618muL}$ (see Corollary 10) yields $D=\exists r.D_{1}\sqcap\ldots\sqcap\exists r.D_{n}$ for $n\geq 0$ $\mathcal{E\mkern-1.618muL}$ concept descriptions $D_{1},\ldots,D_{n}$ satisfying $C\sqsubseteq D_{1},\ldots,C\sqsubseteq D_{n}$ (see Lemma 11). Thus, if we define $\lambda:=\{z_{1}\mapsto D_{1},\ldots,z_{n}\mapsto D_{n},z\mapsto C\}$, then $\lambda\sigma_{n}\approx_{\mathcal{E\mkern-1.618muL}}^{\{x,y\}}\sigma$, which shows $\sigma_{n}\leq_{\mathcal{E\mkern-1.618muL}}^{\{x,y\}}\sigma$. $\hfill\blacktriangleleft$\par
 
% END S3.SS0.SSS0.P0.SPx1.p1
 
% BEGIN Thmtheorem14
 
\par\medskip\noindent\textbf{Lemma 14.}\par
 
% BEGIN Thmtheorem14.p1
 
If $\sigma_{n}$ for $n\geq 0$ are defined as in Example 12, then $\sigma_{0}>_{\mathcal{E\mkern-1.618muL}}^{\{x,y\}}\sigma_{1}>_{\mathcal{E%
\mkern-1.618muL}}^{\{x,y\}}\sigma_{2}>_{\mathcal{E\mkern-1.618muL}}^{\{x,y\}}\cdots$.\par
 
% END Thmtheorem14.p1
 
% END Thmtheorem14
 
% END S3.SS0.SSS0.P0.SPx1
 
% BEGIN S3.SS0.SSS0.P0.SPx2
 
\par\medskip\noindent\textbf{Proof.}\par
 
% BEGIN S3.SS0.SSS0.P0.SPx2.p1
 
First note that $\sigma_{n+1}\leq_{\mathcal{E\mkern-1.618muL}}^{\{x,y\}}\sigma_{n}$ since $\lambda\sigma_{n+1}\approx_{\mathcal{E\mkern-1.618muL}}^{\{x,y\}}\sigma_{n}$ for $\lambda:=\{z_{n+1}\mapsto\top\}$. Second, assume that $\sigma_{n}\leq_{\mathcal{E\mkern-1.618muL}}^{\{x,y\}}\sigma_{n+1}$, i.e., there is a substitution $\tau$ such that $\tau\sigma_{n}\approx_{\mathcal{E\mkern-1.618muL}}^{\{x,y\}}\sigma_{n+1}$. Then $\exists r.\tau(z_{1})\sqcap\ldots\sqcap\exists r.\tau(z_{n})\approx_{\mathcal{%
E\mkern-1.618muL}}\exists r.z_{1}\sqcap\ldots\sqcap\exists r.z_{n+1}$. Consequently, by Theorem 9, the reduced forms of these two concept descriptions would need to be equal up to associativity and commutativity of $\sqcap$. However, the reduced form of the concept description on the left-hand side has at most $n$ existential restrictions, whereas the reduced form of the concept description on the right-hand side has $n+1$ existential restrictions, which yields a contradiction. Thus, we have shown $\sigma_{n+1}<_{\mathcal{E\mkern-1.618muL}}^{\{x,y\}}\sigma_{n}$. $\hfill\blacktriangleleft$\par
 
% END S3.SS0.SSS0.P0.SPx2.p1
 
% BEGIN S3.SS0.SSS0.P0.SPx2.p2
 
The following theorem is now an immediate consequence of Lemma 4.\par
 
% END S3.SS0.SSS0.P0.SPx2.p2
 
% BEGIN Thmtheorem15
 
\par\medskip\noindent\textbf{Theorem 15.}\par
 
% BEGIN Thmtheorem15.p1
 
The restricted unification type of $\mathcal{E\mkern-1.618muL}$ is zero.\par
 
% END Thmtheorem15.p1
 
% END Thmtheorem15
 
% BEGIN S3.SS0.SSS0.P0.SPx2.p3
 
With respect to the unrestricted instantiation preorder, the instance relationship between the substitutions $\sigma_{n}$ and $\sigma_{n+1}$ no longer holds. More generally, we can show the following result.\par
 
% END S3.SS0.SSS0.P0.SPx2.p3
 
% BEGIN Thmtheorem16
 
\par\medskip\noindent\textbf{Lemma 16.}\par
 
% BEGIN Thmtheorem16.p1
 
If $\sigma_{n}$ for $n\geq 0$ are defined as in Example 12, then $\sigma_{n}\not\leq_{\mathcal{E\mkern-1.618muL}}^{V}\sigma_{m}$ for all distinct $n,m\geq 0$.\par
 
% END Thmtheorem16.p1
 
% END Thmtheorem16
 
% END S3.SS0.SSS0.P0.SPx2
 
% BEGIN S3.SS0.SSS0.P0.SPx3
 
\par\medskip\noindent\textbf{Proof.}\par
 
% BEGIN S3.SS0.SSS0.P0.SPx3.p1
 
Let $n<m$. Then we know from Lemma 14 that $\sigma_{n}\not\leq_{\mathcal{E\mkern-1.618muL}}^{\{x,y\}}\sigma_{m}$, which implies $\sigma_{n}\not\leq_{\mathcal{E\mkern-1.618muL}}^{V}\sigma_{m}$ since ${\leq_{\mathcal{E\mkern-1.618muL}}^{V}}\subseteq{\leq_{\mathcal{E\mkern-1.618%
muL}}^{\{x,y\}}}$.\par
 
% END S3.SS0.SSS0.P0.SPx3.p1
 
% BEGIN S3.SS0.SSS0.P0.SPx3.p2
 
Assume that $n>m$ and that there is a substitution $\lambda$ such that $\lambda\sigma_{n}\approx_{\mathcal{E\mkern-1.618muL}}\sigma_{m}$. Then $\lambda\sigma_{n}(x)\approx_{\mathcal{E\mkern-1.618muL}}\sigma_{m}(x)$, and since $\sigma_{m}(x)$ does not contain $z_{n}$, but $\sigma_{n}(x)$ does, $\lambda$ must replace $z_{n}$ by a concept description not containing $z_{n}$, and thus in particular $\lambda(z_{n})\not\approx_{\mathcal{E\mkern-1.618muL}}z_{n}$. But then $\lambda\sigma_{n}(z_{n})=\lambda(z_{n})\not\approx_{\mathcal{E\mkern-1.618muL}%
}z_{n}$ and $\sigma_{m}(z_{n})=z_{n}$ yield a contradiction to $\lambda\sigma_{n}\approx_{\mathcal{E\mkern-1.618muL}}\sigma_{m}$. $\hfill\blacktriangleleft$\par
 
% END S3.SS0.SSS0.P0.SPx3.p2
 
% BEGIN S3.SS0.SSS0.P0.SPx3.p3
 
Thus, w.r.t. the unrestricted instantiation preorder, the set $\{\sigma_{n}\mid n\geq 0\}$ consists of unifiers that are incomparable. However, this set is also no longer complete.\par
 
% END S3.SS0.SSS0.P0.SPx3.p3
 
% BEGIN Thmtheorem17
 
\par\medskip\noindent\textbf{Lemma 17.}\par
 
% BEGIN Thmtheorem17.p1
 
Let $\Gamma$ and $\sigma_{n}$ for $n\geq 0$ be defined as in Example 12. The $\mathcal{E\mkern-1.618muL}$-unifier $\sigma:=\{x\mapsto\exists r.A,y\mapsto A\}$ of $\Gamma$ is not an instance of any of the substitutions $\sigma_{n}$ for $n\geq 0$ if we use the unrestricted instantiation preorder.\par
 
% END Thmtheorem17.p1
 
% END Thmtheorem17
 
% END S3.SS0.SSS0.P0.SPx3
 
% BEGIN S3.SS0.SSS0.P0.SPx4
 
\par\medskip\noindent\textbf{Proof.}\par
 
% BEGIN S3.SS0.SSS0.P0.SPx4.p1
 
First, note that $\sigma_{0}\leq_{\mathcal{E\mkern-1.618muL}}^{V}\sigma$ is not possible since $\sigma_{0}(x)=\top$, and thus $\lambda\sigma_{0}(x)=\top$ holds for all substitutions $\lambda$, whereas $\sigma(x)=\exists r.A$. If $n\geq 1$ and $\lambda\sigma_{n}\approx_{\mathcal{E\mkern-1.618muL}}\sigma$, then $\lambda$ must replace all variables $z_{i}$ occurring in $\sigma_{n}(x)$ with $A$, and thus $\lambda\sigma_{n}(z_{i})=\lambda(z_{i})\approx_{\mathcal{E\mkern-1.618muL}}A$. However, $\sigma(z_{i})=z_{i}$, which yields a contradiction to $\lambda\sigma_{n}\approx_{\mathcal{E\mkern-1.618muL}}\sigma$. $\hfill\blacktriangleleft$\par
 
% END S3.SS0.SSS0.P0.SPx4.p1
 
% BEGIN S3.SS0.SSS0.P0.SPx4.p2
 
Summing up, we have seen that the $\mathcal{E\mkern-1.618muL}$-unification problem $\Gamma$ of Example 12 has unification type zero for the restricted instantiation preorder, but the proof that we have used to show this result does not work if we employ the unrestricted instantiation preorder instead.\par
 
% END S3.SS0.SSS0.P0.SPx4.p2
 
% BEGIN S3.SS0.SSS0.P0.SPx4.p3
 
In the following, we prove the general result that the unrestricted unification type of $\mathcal{E\mkern-1.618muL}$ is not zero. Note that, by Theorem 8, this type cannot be unitary or finitary, and thus must be infinitary. The main idea underlying our proof of this result is to show that, up to the equivalence relation $\sim_{\mathcal{E\mkern-1.618muL}}^{V}$, every substitution has only finitely many more general substitutions w.r.t. $\leq_{\mathcal{E\mkern-1.618muL}}^{V}$. The most challenging technical result needed to prove this is the following lemma.\par
 
% END S3.SS0.SSS0.P0.SPx4.p3
 
% BEGIN Thmtheorem18
 
\par\medskip\noindent\textbf{Lemma 18.}\par
 
% BEGIN Thmtheorem18.p1
 
Let $\sigma,\theta$ be substitutions such that $\sigma\leq_{\mathcal{E\mkern-1.618muL}}^{V}\theta$. Then there is a substitution $\sigma^{\prime}$ such that $\sigma\sim_{\mathcal{E\mkern-1.618muL}}^{V}\sigma^{\prime}$ and $\operatorname{Dom}(\sigma^{\prime})\subseteq\operatorname{Dom}(\theta)$.\par
 
% END Thmtheorem18.p1
 
% END Thmtheorem18
 
% END S3.SS0.SSS0.P0.SPx4
 
% BEGIN S3.SS0.SSS0.P0.SPx5
 
\par\medskip\noindent\textbf{Proof.}\par
 
% BEGIN S3.SS0.SSS0.P0.SPx5.p1
 
First note that we can assume without loss of generality that $\operatorname{VRan}(\sigma)\cap\operatorname{Dom}(\theta)=\emptyset$. Otherwise, we can apply a permutation to $\sigma$ that renames the variables in $\operatorname{VRan}(\sigma)$ appropriately, which yields a substitution that is $\sim_{\mathcal{E\mkern-1.618muL}}^{V}$-equivalent to $\sigma$ and satisfies the required disjointness condition (see the proof of Lemma 5 for how such a permutation can be obtained).\par
 
% END S3.SS0.SSS0.P0.SPx5.p1
 
% BEGIN S3.SS0.SSS0.P0.SPx5.p2
 
Since $\sigma\leq_{\mathcal{E\mkern-1.618muL}}^{V}\theta$, we know that there is a substitution $\lambda$ such that $\lambda\sigma\approx_{\mathcal{E\mkern-1.618muL}}\theta$. Consider the set $\mathcal{X}$ of all variables $x$ such that $x\in\operatorname{Dom}(\sigma)\setminus\operatorname{Dom}(\theta)$. Then $\lambda\sigma(x)\approx_{\mathcal{E\mkern-1.618muL}}\theta(x)=x$ holds for all $x\in\mathcal{X}$. Consequently, the top-level conjunction of $\sigma(x)$ cannot contain concept constants (i.e., constants) or existential restrictions (i.e., terms starting with a function application). This means that $\sigma(x)$ is a conjunction of variables:\par
    \[\sigma(x)=z_{1}^{x}\sqcap\ldots\sqcap z_{n_{x}}^{x},\]    
where we assume (without loss of generality) that all the variables $z_{i}^{x}$ for a fixed $x$ are pairwise distinct. To obtain $\lambda\sigma(x)\approx_{\mathcal{E\mkern-1.618muL}}x$, the substitution $\lambda$ must thus assign (modulo $\approx_{\mathcal{E\mkern-1.618muL}}$) $x$ or $\top$ to the variables $z_{i}^{x}$ for all $i,1\leq i\leq n_{x}$, and $x$ to at least one of these variables. Let us assume without loss of generality that $\lambda$ assigns $x$ to $z_{1}^{x}$ for all $x\in\mathcal{X}$. This implies that $z_{1}^{x}\neq z_{1}^{y}$ for different elements $x,y$ of $\mathcal{X}$.\par
 
% END S3.SS0.SSS0.P0.SPx5.p2
 
% BEGIN S3.SS0.SSS0.P0.SPx5.p3
 
Since $\operatorname{VRan}(\sigma)\cap\operatorname{Dom}(\theta)=\emptyset$, we know that $z_{i}^{x}\not\in\operatorname{Dom}(\theta)$ for all $x\in\mathcal{X}$ and $1\leq i\leq n_{x}$. We claim that $z_{i}^{x}\in\operatorname{Dom}(\sigma)$ also holds, and thus $z_{i}^{x}\in\mathcal{X}$. Otherwise, $\lambda\sigma(z_{i}^{x})=\lambda(z_{i}^{x})\in\{x,\top\}$. However, $\lambda\sigma(z_{i}^{x})\approx_{\mathcal{E\mkern-1.618muL}}\theta(z_{i}^{x})=%
z_{i}^{x}$. This yields a contradiction unless $z_{i}^{x}=x$. But then we also have $z_{i}^{x}\in\operatorname{Dom}(\sigma)$ since $x\in\operatorname{Dom}(\sigma)$. We have thus shown that all the variables $z_{i}^{x}$ for $x\in\mathcal{X}$ also belong to $\mathcal{X}=\operatorname{Dom}(\sigma)\setminus\operatorname{Dom}(\theta)$. Since $z_{1}^{x}\neq z_{1}^{y}$ for different elements $x,y$ of $\mathcal{X}$, this implies that the $z$-variables with index $1$ already “use up” all of $\mathcal{X}$, and thus $n_{x}=1$ holds for all $x\in\mathcal{X}$.\par
 
% END S3.SS0.SSS0.P0.SPx5.p3
 
% BEGIN S3.SS0.SSS0.P0.SPx5.p4
 
Consequently, $\sigma$ is a $W$-renaming for $W=\operatorname{Dom}(\sigma)\setminus\operatorname{Dom}(\theta)$, where according to Definition 2.11 in [28] a substitution $\tau$ is a $W$-renaming if $\tau(x)$ is a variable for all $x\in W$ and $\tau$ is injective on $W$. Since $\sigma$ is the identity on $V\setminus\operatorname{Dom}(\sigma)$ and elements of $\operatorname{Dom}(\sigma)\setminus\operatorname{Dom}(\theta)$ are mapped to $\operatorname{Dom}(\sigma)\setminus\operatorname{Dom}(\theta)$ by $\sigma$, the substitution $\sigma$ is also a $W$-renaming for $W=V\setminus\operatorname{Dom}(\theta)$. By Lemma 2.12 in [28], there is a permutation $\pi$ that coincides with $\sigma$ on $W=V\setminus\operatorname{Dom}(\theta)$. Let $\sigma^{\prime}:=\pi^{-1}\sigma$. Then $\sigma^{\prime}\sim_{\emptyset}^{V}\sigma$, and thus also $\sigma^{\prime}\sim_{\mathcal{E\mkern-1.618muL}}^{V}\sigma$. In addition, for all $x\not\in\operatorname{Dom}(\theta)$ we have $\sigma^{\prime}(x)=\pi^{-1}\sigma(x)=x$, which yields $x\not\in\operatorname{Dom}(\sigma^{\prime})$. This shows that $\operatorname{Dom}(\sigma^{\prime})\subseteq\operatorname{Dom}(\theta)$. $\hfill\blacktriangleleft$\par
 
% END S3.SS0.SSS0.P0.SPx5.p4
 
% BEGIN S3.SS0.SSS0.P0.SPx5.p5
 
This lemma together with the next result shows that, up to equivalence, we can bound the variables occurring in more general substitutions by the substitutions they have as instance.\par
 
% END S3.SS0.SSS0.P0.SPx5.p5
 
% BEGIN Thmtheorem19
 
\par\medskip\noindent\textbf{Lemma 19.}\par
 
% BEGIN Thmtheorem19.p1
 
Let $\sigma,\theta$ be substitutions such that $\sigma\leq_{\mathcal{E\mkern-1.618muL}}^{V}\theta$ and $\operatorname{Dom}(\sigma)\subseteq\operatorname{Dom}(\theta)$. Then $\operatorname{VRan}(\sigma)\subseteq\operatorname{Dom}(\theta)\cup%
\operatorname{VRan}(\theta)$.\par
 
% END Thmtheorem19.p1
 
% END Thmtheorem19
 
% END S3.SS0.SSS0.P0.SPx5
 
% BEGIN S3.SS0.SSS0.P0.SPx6
 
\par\medskip\noindent\textbf{Proof.}\par
 
% BEGIN S3.SS0.SSS0.P0.SPx6.p1
 
Let $\lambda$ be such that $\lambda\sigma\approx_{\mathcal{E\mkern-1.618muL}}\theta$. Assume that $z\in\operatorname{VRan}(\sigma)$, but $z\notin\operatorname{Dom}(\theta)\cup\operatorname{VRan}(\theta)$. Then $z\not\in\operatorname{Dom}(\sigma)$, and thus $\lambda(z)=\lambda\sigma(z)\approx_{\mathcal{E\mkern-1.618muL}}\theta(z)=z$. However, since $z\in\operatorname{VRan}(\sigma)$, there is a variable $x\in\operatorname{Dom}(\sigma)$ such that $\sigma(x)$ contains $z$. Then $x\in\operatorname{Dom}(\theta)$, and thus all variables occurring in $\theta(x)$ belong to $\operatorname{VRan}(\theta)$. Since $\lambda\sigma(x)\approx_{\mathcal{E\mkern-1.618muL}}\theta(x)$, this means that $z\in\operatorname{VRan}(\theta)$, which contradicts our assumptions on $z$. $\hfill\blacktriangleleft$\par
 
% END S3.SS0.SSS0.P0.SPx6.p1
 
% BEGIN S3.SS0.SSS0.P0.SPx6.p2
 
The following lemma bounds the role depth (i.e., the maximal nesting of existential restrictions) as well as the concept constants and role names occurring in more general substitutions. Formally, the role depth $\operatorname{rd}(C)$ of an $\mathcal{E\mkern-1.618muL}$ concept description $C$ is defined inductively as follows:\par
  
\par\noindent  $\blacksquare$ 
% BEGIN S3.I1.i1.p1
 
$\operatorname{rd}(A):=\operatorname{rd}(\top):=0$ for all concept names $A$,\par
 
% END S3.I1.i1.p1
 
 
\par\noindent  $\blacksquare$ 
% BEGIN S3.I1.i2.p1
 
$\operatorname{rd}(C\sqcap D):=\max\{\operatorname{rd}(C),\operatorname{rd}(D)\}$ and $\operatorname{rd}(\exists r.C):=1+\operatorname{rd}(C)$.\par
 
% END S3.I1.i2.p1
 
  
% END S3.SS0.SSS0.P0.SPx6.p2
 
% BEGIN Thmtheorem20
 
\par\medskip\noindent\textbf{Lemma 20.}\par
 
% BEGIN Thmtheorem20.p1
 
Let $\sigma,\theta$ be substitutions such that $\sigma\leq_{\mathcal{E\mkern-1.618muL}}^{V}\theta$. Then the following holds for all $x\in V$: the role depth of $\sigma(x)$ is bounded by the role depth of $\theta(x)$, and the concept constants and role names occurring in $\sigma(x)$ also occur in $\theta(x)$.\par
 
% END Thmtheorem20.p1
 
% END Thmtheorem20
 
% END S3.SS0.SSS0.P0.SPx6
 
% BEGIN S3.SS0.SSS0.P0.SPx7
 
\par\medskip\noindent\textbf{Proof.}\par
 
% BEGIN S3.SS0.SSS0.P0.SPx7.p1
 
This is an easy consequence of the fact that $\approx_{\mathcal{E\mkern-1.618muL}}$ preserves the role depth as well as the set of concept constants and role names occurring in a concept description, and applying a substitution to a concept description can only increase the role depth and add concept constants or role names, but not decrease the role depth or remove concept or role names. $\hfill\blacktriangleleft$\par
 
% END S3.SS0.SSS0.P0.SPx7.p1
 
% BEGIN S3.SS0.SSS0.P0.SPx7.p2
 
As a consequence of the sequence of lemmas we have just shown we know that, for a given substitution $\theta$, the set of more general substitutions is finite up to the equivalence relation $\sim_{\mathcal{E\mkern-1.618muL}}^{V}$. In fact, Lemmas 18 and 19 show that one can restrict the attention to substitutions $\sigma$ satisfying $\operatorname{Dom}(\sigma)\cup\operatorname{VRan}(\sigma)\subseteq%
\operatorname{Dom}(\theta)\cup\operatorname{VRan}(\theta)$. In addition, for all $x\in\operatorname{Dom}(\sigma)$, the concept descriptions $\sigma(x)$ are built using role names and concept constants occurring in $\theta(x)$ as well as variables from $\operatorname{Dom}(\theta)\cup\operatorname{VRan}(\theta)$, and have a role depth that is bounded by the one of $\theta(x)$. It is well-known [13] that there are up to equivalence $\approx_{\mathcal{E\mkern-1.618muL}}$ only finitely many $\mathcal{E\mkern-1.618muL}$-concept descriptions satisfying these properties.\par
 
% END S3.SS0.SSS0.P0.SPx7.p2
 
% BEGIN Thmtheorem21
 
\par\medskip\noindent\textbf{Lemma 21.}\par
 
% BEGIN Thmtheorem21.p1
 
For a given substitution $\theta$, the set $\{\sigma\mid\sigma\leq_{\mathcal{E\mkern-1.618muL}}^{V}\theta\}$ is finite up to $\sim_{\mathcal{E\mkern-1.618muL}}^{V}$-equivalence.\par
 
% END Thmtheorem21.p1
 
% END Thmtheorem21
 
% BEGIN S3.SS0.SSS0.P0.SPx7.p3
 
As an immediate consequence we obtain that the set of substitutions more general than $\theta$ contains elements that are minimal w.r.t. $\leq_{\mathcal{E\mkern-1.618muL}}^{V}$. This obviously shows that the unrestricted unification type of $\mathcal{E\mkern-1.618muL}$ cannot be zero. Since the restricted unification type is zero, the unrestricted type cannot be unitary or finitary by Theorem 8.\par
 
% END S3.SS0.SSS0.P0.SPx7.p3
 
% BEGIN Thmtheorem22
 
\par\medskip\noindent\textbf{Theorem 22.}\par
 
% BEGIN Thmtheorem22.p1
 
The unrestricted unification type of $\mathcal{E\mkern-1.618muL}$ is infinitary.\par
 
% END Thmtheorem22.p1
 
% END Thmtheorem22
 
% END S3.SS0.SSS0.P0.SPx7
 
% END S3


% BEGIN S4
 
\par\medskip\noindent\textbf{4 The unrestricted unification types of ACUI, ACU, and AC}\par
 
% BEGIN S4.p1
 
We consider a signature consisting of a binary function symbol $f$ and a constant symbol $0$. The theory $\mathsf{ACU}$ consists of identities stating that $f$ is associative and commutative and that $0$ is a unit for $f$. The theory $\mathsf{ACUI}$ additionally states that $f$ is idempotent. The theory $\mathsf{AC}$ is obtained from $\mathsf{ACU}$ by removing the unit $0$ from the signature and the identity involving it from the axiomatization. For elementary unification (where unification problems may only contain $f$, $0$, and variables), $\mathsf{ACU}$ and $\mathsf{ACUI}$ are unitary w.r.t. the restricted instantiation preorder, whereas $\mathsf{AC}$ is finitary [19]. In the following, we show that all three theories are infinitary w.r.t. the unrestricted instantiation preorder. These results extend the one from [10] in two directions. First, [10] considers only $\mathsf{ACUI}$, whereas here we also investigate $\mathsf{ACU}$ and $\mathsf{AC}$. Second, [10] provides only the lower bound “at least infinitary” for $\mathsf{ACUI}$, whereas here we determine the exact unification type (infinitary) for the three theories.\par
 
% END S4.p1
 
% BEGIN S4.p2
 
We start with proving that the unrestricted unification type of $\mathsf{AC}$, $\mathsf{ACU}$, and $\mathsf{ACUI}$ is at least infinitary. In contrast to $\mathcal{E\mkern-1.618muL}$, where we were able to use Theorem 8 to deduce that it is not unitary or finitary w.r.t. the unrestricted instantiation preorder from the fact that it is of type zero for the restricted instantiation preorder, we must prove this directly for $\mathsf{ACUI}$, $\mathsf{ACU}$, and $\mathsf{AC}$. Actually, following [10], we show a more general result that holds for regular theories satisfying certain additional restrictions.\par
 
% END S4.p2
 
% BEGIN S4.p3
 
Recall that a set of identities $E$ is regular if $\operatorname{Var}(s)=\operatorname{Var}(t)$ holds for all identities $s\approx t$ in $E$. Regularity of the defining set of identities of an equational theory implies regularity of the whole theory, i.e., if $E$ is a regular set of identities, then $\operatorname{Var}(s)=\operatorname{Var}(t)$ holds for all terms $s,t$ satisfying $s\approx_{E}t$ [58]. The identities of $\mathsf{AC}$, $\mathsf{ACU}$, and $\mathsf{ACUI}$ are regular, and thus the following result from [10] applies to them.\par
 
% END S4.p3
 
% BEGIN Thmtheorem23
 
\par\medskip\noindent\textbf{Lemma 23 ([10]).}\par
 
% BEGIN Thmtheorem23.p1
 
Let $E$ be a regular theory and $\Gamma=\{s\approx_{E}^{?}t\}$ an $E$-unification problem s.t. $\operatorname{Var}(s)\cap\operatorname{Var}(t)=\emptyset$. Then the set $\mathcal{C}_{E}(\Gamma)$ consisting of all $E$-unifiers $\sigma$ of $\Gamma$ satisfying\par
    \[\forall y\in\operatorname{VRan}(\sigma).\exists x,x^{\prime}\in V\text{ s.t. }%
x\neq x^{\prime}\text{ and }y\in\operatorname{Var}(\sigma(x))\cap\operatorname%
{Var}(\sigma(x^{\prime}))\]    
is complete w.r.t. $\leq_{E}^{V}$.\par
 
% END Thmtheorem23.p1
 
% END Thmtheorem23
 
% BEGIN S4.p4
 
Together with Theorem 3, this lemma yields the following result.\par
 
% END S4.p4
 
% BEGIN Thmtheorem24
 
\par\medskip\noindent\textbf{Lemma 24.}\par
 
% BEGIN Thmtheorem24.p1
 
Let $E$ be a regular theory and $\Gamma=\{s\approx_{E}^{?}t\}$ an $E$-unification problem s.t. $\operatorname{Var}(s)\cap\operatorname{Var}(t)=\emptyset$. If $\Gamma$ has a minimal complete set of $E$-unifiers w.r.t. $\leq_{E}^{V}$, then it has one that is contained in $\mathcal{C}_{E}(\Gamma)$.\par
 
% END Thmtheorem24.p1
 
% END Thmtheorem24
 
% BEGIN S4.SS0.SSS0.P0.SPx1
 
\par\medskip\noindent\textbf{Proof.}\par
 
% BEGIN S4.SS0.SSS0.P0.SPx1.p1
 
Since $\Gamma$ has a minimal complete set of $E$-unifiers, the set $M$ of minimal elements of $[\mathcal{U}_{E}(\Gamma)]_{E}^{V}$ is complete. For $\mathcal{C}_{E}(\Gamma)$ to be complete, it must contain for every equivalence class in $M$ at least one representative. Thus, by selecting for each class in $M$ one of its representatives in $\mathcal{C}_{E}(\Gamma)$, we obtain a minimal complete set that is contained in $\mathcal{C}_{E}(\Gamma)$. $\hfill\blacktriangleleft$\par
 
% END S4.SS0.SSS0.P0.SPx1.p1
 
% BEGIN S4.SS0.SSS0.P0.SPx1.p2
 
We are now ready to formulate the “additional restrictions” mentioned above.\par
 
% END S4.SS0.SSS0.P0.SPx1.p2
 
% BEGIN Thmtheorem25
 
\par\medskip\noindent\textbf{Definition 25.}\par
 
% BEGIN Thmtheorem25.p1
 
Given a regular equational theory $E$, we say that the $E$-unification problem $\Gamma$ is NUOF if the following conditions are satisfied:\par
  
\par\noindent  $\blacksquare$ 
% BEGIN S4.I1.i1.p1
 
$\Gamma=\{s\approx_{E}^{?}t\}$ for terms $s,t$ satisfying $\operatorname{Var}(s)\cap\operatorname{Var}(t)=\emptyset$,\par
 
% END S4.I1.i1.p1
 
 
\par\noindent  $\blacksquare$ 
% BEGIN S4.I1.i2.p1
 
there is a $\leq_{E}^{V}$-minimal unifier $\sigma$ of $\Gamma$ that uses fresh variables, i.e., $\operatorname{VRan}(\sigma)\setminus X\neq\emptyset$ where $X=\operatorname{Var}(s)\cup\operatorname{Var}(t)$, and\par
 
% END S4.I1.i2.p1
 
 
\par\noindent  $\blacksquare$ 
% BEGIN S4.I1.i3.p1
 
this unifier $\sigma$ belongs to the set $\mathcal{C}_{E}(\Gamma)$ defined in the formulation of Lemma 23.\par
 
% END S4.I1.i3.p1
 
  
% END Thmtheorem25.p1
 
% END Thmtheorem25
 
% BEGIN S4.SS0.SSS0.P0.SPx1.p3
 
Intuitively, NUOF stands for “not unitary or finitary,” but we still need to show that this name is justified. Given a NUOF $E$-unification problem $\Gamma$, let $x_{0}\in\operatorname{VRan}(\sigma)\setminus X$ and consider the following construction of substitutions:\par
    \[\sigma_{z}:=\sigma\tau_{x_{0},z}\ \ \mbox{where $z\in{V}$ and $\tau_{x_{0},z}:%
=\{x_{0}\mapsto z,z\mapsto x_{0}\}$}.\]    
One can show that, under certain conditions on $z$, such substitutions $\sigma_{z}$ are $\leq_{E}^{V}$-minimal unifiers that are incomparable to each other w.r.t. $\leq_{E}^{V}$. By Theorem 3, this implies that $\Gamma$ cannot have a finite minimal complete set of unifiers w.r.t. $\leq_{E}^{V}$ since there are infinitely many variables $z$ satisfying these conditions.\par
 
% END S4.SS0.SSS0.P0.SPx1.p3
 
% BEGIN Thmtheorem26
 
\par\medskip\noindent\textbf{Lemma 26 ([10]).}\par
 
% BEGIN Thmtheorem26.p1
 
Let $E$ be a regular equational theory $E$, $\Gamma$ a NUOF $E$-unification problem, and $X$ and $\sigma_{z}$ for $z\in V$ be defined as above.\par
  
\par\noindent  $\blacksquare$ 
% BEGIN S4.I2.i1.p1
 
For each $z\in V\setminus X$, $\sigma_{z}$ is a minimal $E$-unifier of $\Gamma$ w.r.t. $\leq_{E}^{V}$.\par
 
% END S4.I2.i1.p1
 
 
\par\noindent  $\blacksquare$ 
% BEGIN S4.I2.i2.p1
 
For any two distinct variables $z,z^{\prime}\in V\setminus(\operatorname{Dom}(\sigma)\cup\operatorname{VRan}(%
\sigma))$, $\sigma_{z}$ and $\sigma_{z^{\prime}}$ are incomparable w.r.t. $\leq_{E}^{V}$.\par
 
% END S4.I2.i2.p1
 
  
% END Thmtheorem26.p1
 
% END Thmtheorem26
 
% BEGIN S4.SS0.SSS0.P0.SPx1.p4
 
Since $V\setminus(X\cup\operatorname{Dom}(\sigma)\cup\operatorname{VRan}(\sigma))$ is infinite, this lemma together with Theorem 3 implies that $\Gamma$ cannot have a finite minimal complete set w.r.t. $\leq_{E}^{V}$.\par
 
% END S4.SS0.SSS0.P0.SPx1.p4
 
% BEGIN Thmtheorem27
 
\par\medskip\noindent\textbf{Theorem 27.}\par
 
% BEGIN Thmtheorem27.p1
 
If $E$ is a regular equational theory and $\Gamma$ a NUOF $E$-unification problem, then $\Gamma$ does not have a finite minimal complete set of $E$-unifiers w.r.t. the unrestricted instantiation preorder $\leq_{E}^{V}$.\par
 
% END Thmtheorem27.p1
 
% END Thmtheorem27
 
% BEGIN S4.SS0.SSS0.P0.SPx1.p5
 
We are now ready to apply this result to $\mathsf{ACUI}$.\par
 
% END S4.SS0.SSS0.P0.SPx1.p5
 
% BEGIN Thmtheorem28
 
\par\medskip\noindent\textbf{Corollary 28 ([10]).}\par
 
% BEGIN Thmtheorem28.p1
 
The unrestricted unification type of $\mathsf{ACUI}$ for elementary unification is at least infinitary.\par
 
% END Thmtheorem28.p1
 
% END Thmtheorem28
 
% END S4.SS0.SSS0.P0.SPx1
 
% BEGIN S4.SS0.SSS0.P0.SPx2
 
\par\medskip\noindent\textbf{Proof.}\par
 
% BEGIN S4.SS0.SSS0.P0.SPx2.p1
 
Since $\mathsf{ACUI}$ is regular, it is sufficient to show that there is an $\mathsf{ACUI}$-unification problem $\Gamma$ that is NUOF. According to Corollary 3.6 in [7], any most general unifier (w.r.t. restricted instantiation) of the $\mathsf{ACUI}$-unification problem $\Gamma=\{f(x,f(y,z))\approx_{\mathsf{ACUI}}^{?}f(u,v)\}$ must use a fresh variable. Let $\theta$ be such an mgu.\par
 
% END S4.SS0.SSS0.P0.SPx2.p1
 
% BEGIN S4.SS0.SSS0.P0.SPx2.p2
 
If $\Gamma$ does not have a minimal complete set of $\mathsf{ACUI}$-unifiers w.r.t. unrestricted instantiation, then we are done. Thus, assume that $\Gamma$ has a minimal complete set $\mathcal{S}$ w.r.t. unrestricted instantiation. By Lemma 24, we can assume without loss of generality that $\mathcal{S}\subseteq\mathcal{C}_{E}(\Gamma)$, and by Theorem 3 we know that the elements of $\mathcal{S}$ are $\leq_{\mathsf{ACUI}}^{V}$-minimal. Since $\theta$ is an $\mathsf{ACUI}$-unifier of $\Gamma$, there is a $\sigma\in\mathcal{S}$ such that $\sigma\leq_{\mathsf{ACUI}}^{V}\theta$. Since ${\leq_{\mathsf{ACUI}}^{V}}\subseteq{\leq_{\mathsf{ACUI}}^{\operatorname{Var}(%
\Gamma)}}$, this implies that $\sigma$ is also an mgu of $\Gamma$ w.r.t. the restricted instantiation preorder, and thus it introduces a fresh variable.\par
 
% END S4.SS0.SSS0.P0.SPx2.p2
 
% BEGIN S4.SS0.SSS0.P0.SPx2.p3
 
Consequently, we have shown that $\Gamma$ is NUOF, and thus Theorem 27 is applicable, which proves the corollary. $\hfill\blacktriangleleft$\par
 
% END S4.SS0.SSS0.P0.SPx2.p3
 
% BEGIN S4.SS0.SSS0.P0.SPx2.p4
 
It remains to show that type zero is not possible for $\mathsf{ACUI}$. This is actually an easy consequence of the results for $\mathcal{E\mkern-1.618muL}$ we have shown in the previous section. In fact, for elementary unification in $\mathsf{ACUI}$, we consider a term set that is a subset of the one for $\mathcal{E\mkern-1.618muL}$ if we use the conjunction operator of $\mathcal{E\mkern-1.618muL}$ as $f$ and the top concept of $\mathcal{E\mkern-1.618muL}$ as unit $0$. On such terms (which we will call $\mathsf{ACUI}$-terms in the following), the equational theories $\approx_{\mathcal{E\mkern-1.618muL}}$ and $\approx_{\mathsf{ACUI}}$ coincide. If we consider substitutions $\sigma,\theta$ using only $\mathsf{ACUI}$-terms, then $\sigma\leq_{\mathsf{ACUI}}^{V}\theta$ clearly implies $\sigma\leq_{\mathcal{E\mkern-1.618muL}}^{V}\theta$. To show the other direction, assume that $\lambda\sigma\approx_{\mathcal{E\mkern-1.618muL}}\theta$, but there is a variable $z$ such that $\lambda(z)$ is not an $\mathsf{ACUI}$ term, i.e., contains a concept constant or an existential restriction. If $z\not\in\operatorname{Dom}(\sigma)$ or $z\in\operatorname{VRan}(\sigma)$, then this leads to a contradiction. Otherwise, we can modify $\lambda$ to $\lambda^{\prime}$ by setting $\lambda^{\prime}(z)=z$ and still have $\lambda^{\prime}\sigma\approx_{\mathcal{E\mkern-1.618muL}}\theta$. This shows that we can assume without loss of generality that $\lambda$ uses only $\mathsf{ACUI}$-terms, which completes the proof that $\sigma\leq_{\mathcal{E\mkern-1.618muL}}^{V}\theta$ implies $\sigma\leq_{\mathsf{ACUI}}^{V}\theta$. Thus, Lemma 21 entails that the set $\{\sigma\mid\sigma\leq_{\mathsf{ACUI}}^{V}\theta\}$ is finite up to $\sim_{\mathcal{E\mkern-1.618muL}}^{V}$-equivalence, and thus also up to $\sim_{\mathsf{ACUI}}^{V}$-equivalence.\par
 
% END S4.SS0.SSS0.P0.SPx2.p4
 
% BEGIN Thmtheorem29
 
\par\medskip\noindent\textbf{Theorem 29.}\par
 
% BEGIN Thmtheorem29.p1
 
The unrestricted unification type of $\mathsf{ACUI}$ for elementary unification is infinitary.\par
 
% END Thmtheorem29.p1
 
% END Thmtheorem29
 
% BEGIN S4.SS0.SSS0.P0.SPx2.p5
 
Theorem 27 also applies to elementary unification in $\mathsf{ACU}$. It is well-known (see, e.g., Section 10.3 in [16]) that a given elementary $\mathsf{ACU}$-unification problem $\Gamma$ can be translated into a system of homogeneous linear diophantine equations. W.r.t. the restricted instantiation preorder, the mgu of the problem $\Gamma$ can then be obtained from the minimal generating set of the solutions of this system, also called its Hilbert base, where the number of variables used in the range of this mgu corresponds to the cardinality of the Hilbert base of the system. As pointed out in Example 2 of [40], the cardinality of the Hilbert base for equations of the form $ny=x_{1}+2x_{2}+\ldots nx_{n}$ grows at least exponentially in $n$, and thus there are clearly instances where the mgu of the corresponding $\mathsf{ACU}$-unification problem needs more than $n+1$ variables. Given this, one can now proceed as in the case of $\mathsf{ACUI}$ to show that Theorem 27 applies.\par
 
% END S4.SS0.SSS0.P0.SPx2.p5
 
% BEGIN Thmtheorem30
 
\par\medskip\noindent\textbf{Corollary 30.}\par
 
% BEGIN Thmtheorem30.p1
 
The unrestricted unification type of $\mathsf{ACU}$ for elementary unification is at least infinitary.\par
 
% END Thmtheorem30.p1
 
% END Thmtheorem30
 
% BEGIN S4.SS0.SSS0.P0.SPx2.p6
 
Proving that type infinitary is also the upper bound in these cases turns out to be more involved than for $\mathsf{ACUI}$. In fact, our proof of the finiteness result used for $\mathcal{E\mkern-1.618muL}$ and then adapted to $\mathsf{ACUI}$ depends on idempotency, and thus does not apply to $\mathsf{ACU}$. Instead, we show that a set of the form $\{\sigma\mid\sigma\leq_{\mathsf{ACU}}^{V}\theta\}$ for a given unifier $\theta$ cannot contain an infinite decreasing chain w.r.t. $>_{\mathsf{ACU}}^{V}$, and thus $\theta$ must be above a minimal unifier. Our proof of this result in Appendix A uses the fact that analogues of Lemmas 18 and 19 also hold for $\mathsf{ACU}$, but requires quite some additional effort to prove the non-existence of infinite decreasing chains.\par
 
% END S4.SS0.SSS0.P0.SPx2.p6
 
% BEGIN Thmtheorem31
 
\par\medskip\noindent\textbf{Theorem 31.}\par
 
% BEGIN Thmtheorem31.p1
 
The unrestricted unification type of $\mathsf{ACU}$ for elementary unification is infinitary.\par
 
% END Thmtheorem31.p1
 
% END Thmtheorem31
 
% BEGIN S4.SS0.SSS0.P0.SPx2.p7
 
The theory $\mathsf{AC}$, which is obtained from $\mathsf{ACU}$ by removing the unit $0$ from the signature and the identity containing it from the axiomatization, is finitary w.r.t. the restricted instantiation preorder. Proving that the unrestricted unification type of $\mathsf{AC}$ is infinitary turns out to be easier than for $\mathsf{ACU}$.\par
 
% END S4.SS0.SSS0.P0.SPx2.p7
 
% BEGIN S4.SS0.SSS0.P0.SPx2.p8
 
Showing that the unrestricted unification type of $\mathsf{AC}$ cannot be unitary or finitary is very similar to our proofs for $\mathsf{ACUI}$ and $\mathsf{ACU}$. In contrast to $\mathsf{ACU}$, the theory $\mathsf{AC}$ is finitary rather than unitary in the restricted setting. A minimal complete set of unifiers is obtained by taking appropriate subsets of the Hilbert base and turning them into unifiers that have a variable in the range for each element of the subset (see, e.g., [16, 32]). Since the full Hilbert base is an appropriate subset, the minimal complete set of $\mathsf{AC}$-unifiers for the unification problem $\Gamma_{n}$ corresponding to the linear diophantine equation $ny=x_{1}+2x_{2}+\ldots nx_{n}$ for a large enough natural number $n$ must contain a unifier $\theta$ that introduces a fresh variable. As in the case of $\mathsf{ACUI}$ we can now show that, under the assumption that $\Gamma$ has a minimal complete set $\mathcal{S}$ of $\mathsf{AC}$-unifiers w.r.t. the unrestricted instantiation preorder, this set $\mathcal{S}$ contains a $\leq_{\mathsf{AC}}^{V}$-minimal $\mathsf{AC}$-unifier $\sigma$ satisfying $\sigma\leq_{\mathsf{AC}}^{V}\theta$, and thus also $\sigma\leq_{\mathsf{AC}}^{\operatorname{Var}(\Gamma_{n})}\theta$. Since $\theta$ contains a variable in the range for each element of the Hilbert base and elements of the Hilbert base cannot be generated by a sum of other vectors, this implies that $\sigma$ must also contain at least as many variables in its range as $\theta$. Consequently, we have shown that $\Gamma_{n}$ is NUOF, and thus Theorem 27 is applicable.\par
 
% END S4.SS0.SSS0.P0.SPx2.p8
 
% BEGIN Thmtheorem32
 
\par\medskip\noindent\textbf{Corollary 32.}\par
 
% BEGIN Thmtheorem32.p1
 
The unrestricted unification type of $\mathsf{AC}$ for elementary unification is at least infinitary.\par
 
% END Thmtheorem32.p1
 
% END Thmtheorem32
 
% BEGIN S4.SS0.SSS0.P0.SPx2.p9
 
To prove that the unrestricted unification type of $\mathsf{AC}$ cannot be zero, we employ the same approach as for unification in $\mathcal{E\mkern-1.618muL}$. First note that Lemmas 18 and 19 also hold if we replace $\mathcal{E\mkern-1.618muL}$ with $\mathsf{AC}$. The main difference is that the crucial argument in the proof of Lemma 18 becomes simpler. There, we consider a setting where $\lambda\sigma(x)\approx_{\mathcal{E\mkern-1.618muL}}\theta(x)=x$ holds for all variables $x\in\mathcal{X}$, and conclude that then $\sigma(x)=z_{1}^{x}\sqcap\ldots\sqcap z_{n_{x}}^{x}$ for some variables $z_{1}^{x},\ldots,z_{n_{x}}^{x}$. Thus is trivially the case for $\mathsf{AC}$ if we replace $\approx_{\mathcal{E\mkern-1.618muL}}$ with $\approx_{\mathsf{AC}}$ and the binary conjunction operator $\sqcap$ of $\mathcal{E\mkern-1.618muL}$ with the $\mathsf{AC}$ function symbol $f$. Then it is proved (by a somewhat involved counting argument) that $n_{x}=1$ for all $x\in\mathcal{X}$ and $z_{1}^{x}\neq z_{1}^{y}$ for distinct elements $x,y\in\mathcal{X}$. For $\mathsf{AC}$, this is again trivially the case since there is no idempotency and the unit $0$ (corresponding to $\top$ in $\mathcal{E\mkern-1.618muL}$) is not available. The rest of the proof of Lemma 18 and the proof of Lemma 19 then work as in the case of $\mathcal{E\mkern-1.618muL}$. Lemma 20 can be replaced by the following result.\par
 
% END S4.SS0.SSS0.P0.SPx2.p9
 
% BEGIN Thmtheorem33
 
\par\medskip\noindent\textbf{Lemma 33.}\par
 
% BEGIN Thmtheorem33.p1
 
Let $\sigma,\theta$ be substitutions such that $\sigma\leq_{\mathsf{AC}}^{V}\theta$. Then the following holds for all $x\in V$: if $\sigma(x)$ and $\theta(x)$ respectively contain $m$ and $n$ occurrences of (not necessarily distinct) variables, then $m\leq n$.\par
 
% END Thmtheorem33.p1
 
% END Thmtheorem33
 
% END S4.SS0.SSS0.P0.SPx2
 
% BEGIN S4.SS0.SSS0.P0.SPx3
 
\par\medskip\noindent\textbf{Proof.}\par
 
% BEGIN S4.SS0.SSS0.P0.SPx3.p1
 
Since we consider elementary unification in $\mathsf{AC}$, all terms are built using the $\mathsf{AC}$ symbol $f$ and variables. If $\lambda\sigma(x)\approx_{\mathsf{AC}}\theta(x)$, then every occurrence of a variable in $\sigma(x)$ is replaced by a term containing one or more variables. Consequently, $\lambda\sigma(x)$ is a term containing occurrences of at least as many variables as $\sigma(x)$. Since the theory cannot change the number of variable occurrences, but only rearrange parentheses and reorder variables, the same is true for $\theta(x)$. $\hfill\blacktriangleleft$\par
 
% END S4.SS0.SSS0.P0.SPx3.p1
 
% BEGIN S4.SS0.SSS0.P0.SPx3.p2
 
Together with the $\mathsf{AC}$ analogues of Lemmas 18 and 19, this lemma yields an $\mathsf{AC}$ analogue of Lemma 21, and thus the following theorem.\par
 
% END S4.SS0.SSS0.P0.SPx3.p2
 
% BEGIN Thmtheorem34
 
\par\medskip\noindent\textbf{Theorem 34.}\par
 
% BEGIN Thmtheorem34.p1
 
The unrestricted unification type of $\mathsf{AC}$ for elementary unification is infinitary.\par
 
% END Thmtheorem34.p1
 
% END Thmtheorem34
 
% BEGIN S4.SS0.SSS0.P0.SPx3.p3
 
Note that, due to the availability of the unit $0$, an $\mathsf{ACU}$-analogue of Lemma 33 does not hold. This is why we need the more involved argument in Appendix A.\par
 
% END S4.SS0.SSS0.P0.SPx3.p3
 
% END S4.SS0.SSS0.P0.SPx3
 
% END S4


% BEGIN A1
 
\par\medskip\noindent\textbf{Appendix A The unification type of ACU}\par
 
% BEGIN A1.p1
 
In this section, we complete the proof of Theorem 31 by showing that type infinitary is also the upper bound for $\mathsf{ACU}$. In contrast to the other equational theories considered in Sections 3 and 4, an analog of Lemma 21 does not hold for $\mathsf{ACU}$. Instead, we show that a set of the form $\{\sigma\mid\sigma\leq_{\mathsf{ACU}}^{V}\theta\}$ for a given unifier $\theta$ cannot contain an infinite decreasing chain w.r.t. $>_{\mathsf{ACU}}^{V}$. We do this in two steps:\par
  
\par\noindent  1. 
% BEGIN A1.I1.i1.p1
 
We recall a well-known characterization of the congruence classes of $\approx_{\mathsf{ACU}}$, and then employ this characterization to show an auxiliary result about $\leq_{\mathsf{ACU}}^{V}$.\par
 
% END A1.I1.i1.p1
 
 
\par\noindent  2. 
% BEGIN A1.I1.i2.p1
 
Finally, we use this auxiliary result to provide the final argument that shows that type zero is not possible for $\mathsf{ACU}$.\par
 
% END A1.I1.i2.p1
 
  
% END A1.p1
 
% BEGIN A1.p2
 
Let $\Sigma=\{f,0\}$ be the signature introduced in Section 4 to define $\mathsf{ACU}$. Given a term $t\in T(\Sigma,V)$, we denote by $c^{t}_{x}$ the number of occurrences of the variable $x$ in $t$. The congruence classes of $\approx_{\mathsf{ACU}}$ can be characterized in the following way (Lemma 10.3.1 of [16]).\par
 
% END A1.p2
 
% BEGIN Thmtheorem35
 
\par\medskip\noindent\textbf{Lemma 35 ([16]).}\par
 
% BEGIN Thmtheorem35.p1
 
Let $s,t\in T(\Sigma,V)$. Then, $s\approx_{\mathsf{ACU}}t$ iff $c^{t}_{x}=c^{s}_{x}$ for all $x\in V$.\par
 
% END Thmtheorem35.p1
 
% END Thmtheorem35
 
% BEGIN A1.p3
 
This lemma tells us that, given a finite set $X_{n}=\{x_{1},\ldots,x_{n}\}\subseteq V$, the $\approx_{\mathsf{ACU}}$-equivalence class of a term $t\in T(\Sigma,X_{n})$ can be uniquely represented by the vector $\vec{v}_{n}(t)=(c^{t}_{x_{1}},\ldots,c^{t}_{x_{n}})$. If $t\approx_{\mathsf{ACU}}0$, then each component $c^{t}_{x_{j}}$ ($1\leq j\leq n$) of $\vec{v}_{n}(t)$ has value zero.\par
 
% END A1.p3
 
% BEGIN A1.p4
 
Based on this representation, a substitution $\sigma$ with $\operatorname{Dom}(\sigma)\subseteq X_{n}$ and $\operatorname{VRan}(\sigma)\subseteq X_{m}$ ($n,m\geq 0$) can be represented as an $n\times m$ matrix with rows $\vec{v}_{m}(\sigma(x_{1})),\ldots,\vec{v}_{m}(\sigma(x_{n}))$. Then, applying a substitution $\sigma$ to a term $t$ corresponds to multiplying the vector for $t$ with the matrix for $\sigma$, i.e., the representation $\vec{v}_{m}(\sigma(t))$ of $\sigma(t)$ can be obtained by computing the following sums (for all $j,1\leq j\leq m$):\par
    \[c^{\sigma(t)}_{x_{j}}\ =\ \sum\limits_{k=1}^{n}c^{t}_{x_{k}}\cdot c^{\sigma(x_%
{k})}_{x_{j}}.\]    
For example, if $t=f(x_{1},f(x_{2},x_{2}))$ and $\sigma=\{x_{1}\mapsto x_{2},x_{2}\mapsto f(x_{1},x_{1})\}$, then $\vec{v}_{2}(t)=(1,2)$ and the matrix representing $\sigma$ has the rows $(0,1)$ and $(2,0)$. The vector representing $\sigma(t)=f(x_{2},f(f(x_{1},x_{1}),f(x_{1},x_{1})))$ is $(4,1)$. Note that $4=1\cdot 0+2\cdot 2$ and $1=1\cdot 1+2\cdot 0$.\par
 
% END A1.p4
 
% BEGIN A1.p5
 
Based on these representations, we next introduce some notions, and show a property about $\leq_{\mathsf{ACU}}^{V}$ that will be useful later on. For simplicity, given a substitution $\sigma$, we will write $(c^{\sigma}_{i1},\ldots,c^{\sigma}_{im})$ instead of $(c^{\sigma(x_{i})}_{x_{1}},\ldots,c^{\sigma(x_{i})}_{x_{m}})$ to refer to the vector $\vec{v}_{m}(\sigma(x_{i}))$. In addition, we denote by $P_{n}$ the set of pairs $\{1,\ldots,n\}\times\{1,\ldots,n\}$.\par
 
% END A1.p5
 
% BEGIN Thmtheorem36
 
\par\medskip\noindent\textbf{Definition 36.}\par
 
% BEGIN Thmtheorem36.p1
 
Let $n\geq 0$ and $P\subseteq P_{n}$. Further, let $\sigma_{1}$ and $\sigma_{2}$ be substitutions whose domains and variable ranges are contained in $X_{n}$. We say that $\sigma_{1}$ is greater than $\sigma_{2}$ w.r.t. $P$ (denoted as $>_{P}$), if\par
  
\par\noindent  $\blacksquare$ 
% BEGIN A1.I2.i1.p1
 
for all $(i,j)\in P$ and $k\in\{1,\ldots,n\}$: $c^{\sigma_{1}}_{ij}\ >\ c^{\sigma_{2}}_{ik}$.\par
 
% END A1.I2.i1.p1
 
  
Further, we say that $\sigma_{1}$ and $\sigma_{2}$ coincide on $P$, if\par
  
\par\noindent  $\blacksquare$ 
% BEGIN A1.I3.i1.p1
 
for all $(i,j)\in P$: $c^{\sigma_{1}}_{ij}\ =c^{\sigma_{2}}_{ij}$.\par
 
% END A1.I3.i1.p1
 
  
% END Thmtheorem36.p1
 
% END Thmtheorem36
 
% BEGIN A1.p6
 
It is not hard to show that $>_{P}$ is a transitive relation.\par
 
% END A1.p6
 
% BEGIN Thmtheorem37
 
\par\medskip\noindent\textbf{Lemma 37.}\par
 
% BEGIN Thmtheorem37.p1
 
Let $n\geq 0$ and $P\subseteq P_{n}$. The relation $>_{P}$ is transitive.\par
 
% END Thmtheorem37.p1
 
% END Thmtheorem37
 
% BEGIN A1.SS0.SSS0.P0.SPx1
 
\par\medskip\noindent\textbf{Proof.}\par
 
% BEGIN A1.SS0.SSS0.P0.SPx1.p1
 
Let $\sigma_{1},\sigma_{2}$ and $\sigma_{3}$ be substitutions such that $\sigma_{1}>_{P}\sigma_{2}$ and $\sigma_{2}>_{P}\sigma_{3}$. To prove transitivity of $>_{P}$, we need to show that $\sigma_{1}>_{P}\sigma_{3}$.\par
 
% END A1.SS0.SSS0.P0.SPx1.p1
 
% BEGIN A1.SS0.SSS0.P0.SPx1.p2
 
Let us take any pair $(i,j)\in P$. Since $\sigma_{1}>_{P}\sigma_{2}$, we have that $c^{\sigma_{1}}_{ij}\ >\ c^{\sigma_{2}}_{ik}$ for all $k\in\{1,\ldots,n\})$. In particular, this yields $c^{\sigma_{1}}_{ij}\ >\ c^{\sigma_{2}}_{ij}$. Furthermore, from $\sigma_{2}>_{P}\sigma_{3}$, we know that $c^{\sigma_{2}}_{ij}\ >\ c^{\sigma_{3}}_{ik}$ for all $k\in\{1,\ldots,n\})$. Consequently, it follows that $c^{\sigma_{1}}_{ij}\ >\ c^{\sigma_{3}}_{ik}\ \ \mbox{for all }k\in\{1,\ldots,n\}.$ Thus, since $(i,j)\in P$ was arbitrarily chosen, we can conclude that $\sigma_{1}>_{P}\sigma_{3}$. $\hfill\blacktriangleleft$\par
 
% END A1.SS0.SSS0.P0.SPx1.p2
 
% BEGIN A1.SS0.SSS0.P0.SPx1.p3
 
We continue by proving the following property about $\leq_{\mathsf{ACU}}^{V}$.\par
 
% END A1.SS0.SSS0.P0.SPx1.p3
 
% BEGIN Thmtheorem38
 
\par\medskip\noindent\textbf{Lemma 38.}\par
 
% BEGIN Thmtheorem38.p1
 
Let $n\geq 0$ and $P\subseteq P_{n}$. In addition, let $\sigma_{1}$, $\sigma_{2},\sigma_{3}$ be substitutions such that:\par
  
\par\noindent  $\blacksquare$ 
% BEGIN A1.I4.i1.p1
 
$\operatorname{Dom}(\theta_{\ell})\cup\operatorname{VRan}(\theta_{\ell})%
\subseteq X_{n}$ for $\ell\in\{1,2,3\}$,\par
 
% END A1.I4.i1.p1
 
 
\par\noindent  $\blacksquare$ 
% BEGIN A1.I4.i2.p1
 
$\sigma_{1}\leq_{\mathsf{ACU}}^{V}\sigma_{2}$,\par
 
% END A1.I4.i2.p1
 
 
\par\noindent  $\blacksquare$ 
% BEGIN A1.I4.i3.p1
 
$\sigma_{1}$ is greater than $\sigma_{2}$ w.r.t. $P$, and\par
 
% END A1.I4.i3.p1
 
 
\par\noindent  $\blacksquare$ 
% BEGIN A1.I4.i4.p1
 
$\sigma_{1}$ and $\sigma_{3}$ coincide on $P_{n}\setminus P$.\par
 
% END A1.I4.i4.p1
 
  
Then, $\sigma_{3}\leq_{\mathsf{ACU}}^{V}\sigma_{2}$.\par
 
% END Thmtheorem38.p1
 
% END Thmtheorem38
 
% END A1.SS0.SSS0.P0.SPx1
 
% BEGIN A1.SS0.SSS0.P0.SPx2
 
\par\medskip\noindent\textbf{Proof.}\par
 
% BEGIN A1.SS0.SSS0.P0.SPx2.p1
 
Since $\sigma_{1}\leq_{\mathsf{ACU}}^{V}\sigma_{2}$, there is a substitution $\lambda$ such that $\lambda\sigma_{1}\approx_{\mathsf{ACU}}\sigma_{2}$. We can assume that the domain and variable range of $\lambda$ satisfy the following properties:\par
  
\par\noindent  $\blacksquare$ 
% BEGIN A1.I5.i1.p1
 
$\operatorname{Dom}(\lambda)\subseteq X_{n}$. For otherwise, if there is a variable $z$ such that $z\in\operatorname{Dom}(\lambda)$ but $z\not\in X_{n}$, then $\operatorname{Dom}(\sigma_{1})\subseteq X_{n}$ implies that $\lambda\sigma_{1}(z)\neq z$. However, this would contradict $\lambda\sigma_{1}\approx_{\mathsf{ACU}}\sigma_{2}$, since $\operatorname{Dom}(\sigma_{2})\subseteq X_{n}$.\par
 
% END A1.I5.i1.p1
 
 
\par\noindent  $\blacksquare$ 
% BEGIN A1.I5.i2.p1
 
$\operatorname{VRan}(\lambda)\subseteq X_{n}$. Note that if $\operatorname{VRan}(\lambda)$ contains a variable $z$ not in $X_{n}$, then $\lambda(x_{i})$ contains $z$ for some $i\in\{1,\ldots,n\}$. But, since $z\not\in\operatorname{VRan}(\sigma_{2})$, the variable $x_{i}$ cannot be in $\operatorname{VRan}(\sigma_{1})$. Thus, we can safely remove the occurrences of $z$ from $\lambda(x_{i})$, i.e., $\lambda\sigma_{1}\approx_{\mathsf{ACU}}\sigma_{2}$ remains true.\par
 
% END A1.I5.i2.p1
 
  
Hence, since $\operatorname{Dom}(\sigma_{1})\cup\operatorname{Dom}(\sigma_{3})\cup%
\operatorname{Dom}(\lambda)\subseteq X_{n}$, we have that $\operatorname{Dom}(\lambda\sigma_{1})\subseteq X_{n}$ and $\operatorname{Dom}(\lambda\sigma_{3})\subseteq X_{n}$. Therefore, to prove that $\sigma_{3}\leq_{\mathsf{ACU}}^{V}\sigma_{2}$, it suffices to show that\par
    \[\lambda\sigma_{1}(x_{i})\ \approx_{\mathsf{ACU}}\ \lambda\sigma_{3}(x_{i})\ \ %
\mbox{for all }i\in\{1,\ldots,n\}.\]  (1)   
Let $i\in\{1,\ldots,n\}$. Since $\operatorname{VRan}(\sigma_{1})\cup\operatorname{VRan}(\lambda)\subseteq X_{n}$, the $\approx_{\mathsf{ACU}}$-equivalence class of $\lambda\sigma_{1}(x_{i})$ can be represented with a vector of the form $(c^{\lambda\sigma_{1}}_{i1},\ldots,c^{\lambda\sigma_{1}}_{in})$, where each value $c^{\lambda\sigma_{1}}_{ij}$ ($1\leq j\leq n$) is determined by the following expression:\par
    \[c^{\lambda\sigma_{1}}_{ij}\ =\ \sum_{k=1}^{n}c^{\sigma_{1}}_{ik}\cdot c^{%
\lambda}_{kj}.\]    
As $\lambda\sigma_{1}(x_{i})\approx_{\mathsf{ACU}}\sigma_{2}(x_{i})$, an application of Lemma 35 yields:\par
    \[c^{\lambda\sigma_{1}}_{ij}\ =\ \sum_{k=1}^{n}c^{\sigma_{1}}_{ik}\cdot c^{%
\lambda}_{kj}\ =\ c^{\sigma_{2}}_{ij}\ \ \mbox{for all }j\in\{1,\ldots,n\}.\]  (2)   
Consider any pair $(i,k)\in P$. Since $\sigma_{1}>_{P}\sigma_{2}$, this means that $c^{\sigma_{1}}_{ik}>c^{\sigma_{2}}_{ij}$ for all $j\in\{1,\ldots,n\}$. Hence, since $\operatorname{VRan}(\lambda)\subseteq X_{n}$, it must be that $\lambda(x_{k})=0$. For otherwise, $c^{\lambda}_{kj}>0$ for some $j\in\{1,\ldots,n\}$ and the expression in (2) would not be true for such $j$. Therefore, (2) can be turned into:\par
    \[c^{\lambda\sigma_{1}}_{ij}\ =\ \sum_{\begin{subarray}{c}k=1\\[2.45836pt]
(i,k)\in P_{n}\setminus P\end{subarray}}^{n}c^{\sigma_{1}}_{ik}\cdot c^{%
\lambda}_{kj}\ =\ c^{\sigma_{2}}_{ij}\ \ \mbox{for all }j\in\{1,\ldots,n\}.\]    
% END A1.SS0.SSS0.P0.SPx2.p1
 
% BEGIN A1.SS0.SSS0.P0.SPx2.p2
 
Regarding $\lambda\sigma_{3}$, we also have $\operatorname{VRan}(\sigma_{3})\cup\operatorname{VRan}(\lambda)\subseteq X_{n}$. Hence, the $\approx_{\mathsf{ACU}}$-equivalence class of $\lambda\sigma_{3}(x_{i})$ has a representation of the form $(c^{\lambda\sigma_{3}}_{i1},\ldots,c^{\lambda\sigma_{3}}_{in})$, where ($1\leq j\leq n$):\par
    \[c^{\lambda\sigma_{3}}_{ij}\ =\ \sum_{\begin{subarray}{c}k=1\\[2.45836pt]
(i,k)\in P_{n}\setminus P\end{subarray}}^{n}c^{\sigma_{3}}_{ik}\cdot c^{%
\lambda}_{kj}.\]    
But then, since $\sigma_{1}$ and $\sigma_{3}$ coincide on $P_{n}\setminus P$, we have that:\par
    \[c^{\lambda\sigma_{3}}_{ij}\ =\ \sum_{\begin{subarray}{c}k=1\\[2.45836pt]
(i,k)\in P_{n}\setminus P\end{subarray}}^{n}c^{\sigma_{3}}_{ik}\cdot c^{%
\lambda}_{kj}\ =\ \sum_{\begin{subarray}{c}k=1\\[2.45836pt]
(i,k)\in P_{n}\setminus P\end{subarray}}^{n}c^{\sigma_{1}}_{ik}\cdot c^{%
\lambda}_{kj}\ =\ c^{\lambda\sigma_{1}}_{ij}.\]    
Hence, we have shown that $(c^{\lambda\sigma_{1}}_{i1},\ldots,c^{\lambda\sigma_{1}}_{in})=(c^{\lambda%
\sigma_{3}}_{i1},\ldots,c^{\lambda\sigma_{3}}_{in})$. Then, an application of Lemma 35 yields that $\lambda\sigma_{1}(x_{i})\approx_{\mathsf{ACU}}\lambda\sigma_{3}(x_{i})$. Thus, since $i$ was arbitrarily selected, we have proved the claim in (1). This concludes the proof of the lemma. $\hfill\blacktriangleleft$\par
 
% END A1.SS0.SSS0.P0.SPx2.p2
 
% BEGIN A1.SS0.SSS0.P0.SPx2.p3
 
We are now ready to move into the final part of our proof. The main idea is to show the following: if $\{\sigma\mid\sigma\leq_{\mathsf{ACU}}^{V}\theta\}$ for a given substitution $\theta$ contains an infinite decreasing chain w.r.t. $>_{\mathsf{ACU}}^{V}$, then such a chain contains three substitutions $\sigma_{1},\sigma_{2}$ and $\sigma_{3}$ such that:\par
  
\par\noindent  $\blacksquare$ 
% BEGIN A1.I6.i1.p1
 
$\sigma_{3}>_{\mathsf{ACU}}^{V}\sigma_{2}>_{\mathsf{ACU}}^{V}\sigma_{1}$, and\par
 
% END A1.I6.i1.p1
 
 
\par\noindent  $\blacksquare$ 
% BEGIN A1.I6.i2.p1
 
$\sigma_{1},\sigma_{2}$ and $\sigma_{3}$ satisfy the hypothesis of Lemma 38.\par
 
% END A1.I6.i2.p1
 
  
This implies that $\sigma_{2}\sim_{\mathsf{ACU}}^{V}\sigma_{3}$, which contradicts the existence of such an infinite chain.\par
 
% END A1.SS0.SSS0.P0.SPx2.p3
 
% BEGIN A1.SS0.SSS0.P0.SPx2.p4
 
In order to apply Lemma 38, we first need to show that there is $n\geq 0$ such that the domain and variable range of each substitution in $\{\sigma\mid\sigma\leq_{\mathsf{ACU}}^{V}\theta\}$ is contained in $X_{n}$. This is a consequence of the analogs of Lemmas 18 and 19 for $\mathsf{ACU}$.\par
 
% END A1.SS0.SSS0.P0.SPx2.p4
 
% BEGIN Thmtheorem39
 
\par\medskip\noindent\textbf{Lemma 39.}\par
 
% BEGIN Thmtheorem39.p1
 
Let $\sigma,\theta$ be substitutions such that $\sigma\leq_{\mathsf{ACU}}^{V}\theta$. Then there is a substitution $\sigma^{\prime}$ such that $\sigma\sim_{\mathsf{ACU}}^{V}\sigma^{\prime}$ and $\operatorname{Dom}(\sigma^{\prime})\subseteq\operatorname{Dom}(\theta)$.\par
 
% END Thmtheorem39.p1
 
% END Thmtheorem39
 
% END A1.SS0.SSS0.P0.SPx2
 
% BEGIN A1.SS0.SSS0.P0.SPx3
 
\par\medskip\noindent\textbf{Proof.}\par
 
% BEGIN A1.SS0.SSS0.P0.SPx3.p1
 
Let $\mathcal{X}=\operatorname{Dom}(\sigma)\setminus\operatorname{Dom}(\theta)$. Since $\sigma\leq_{\mathsf{ACU}}^{V}\theta$, there exists a substitution $\lambda$ such that $\lambda\sigma\approx_{\mathsf{ACU}}\theta$. This means that $\lambda\sigma(x)\approx_{\mathsf{ACU}}\theta(x)=x$ for all $x\in\mathcal{X}$. Hence, an application of Lemma 35 (and consequences thereof described above) yields that:\par
    \[c^{\lambda\sigma(x)}_{x}\ =\ \sum\limits_{z\in V}c^{\sigma(x)}_{z}\cdot c^{%
\lambda(z)}_{x}\ =\ c^{\theta(x)}_{x}\ =\ 1.\]    
This implies that $\sigma(x)$ must contain a single occurrence of some variable $z_{x}$ such that $\lambda(z_{x})=x$, and that $\lambda(y)=0$ for any other variable $y$ occurring in $\sigma(x)$. In addition, $z_{x}\neq z_{y}$ for different elements $x,y$ of $\mathcal{X}$. Once we have this, the same counting argument employed in the proof of Lemma 18 can be applied here to conclude that $\sigma(x)=z_{x}$ for all $x\in\mathcal{X}$.\par
 
% END A1.SS0.SSS0.P0.SPx3.p1
 
% BEGIN A1.SS0.SSS0.P0.SPx3.p2
 
The rest of the proof uses the same arguments as our proof of Lemma 18. $\hfill\blacktriangleleft$\par
 
% END A1.SS0.SSS0.P0.SPx3.p2
 
% BEGIN A1.SS0.SSS0.P0.SPx3.p3
 
The following lemma states that Lemma 19 also holds if we replace $\mathcal{E\mkern-1.618muL}$ with $\mathsf{ACU}$. The proof is the same as for Lemma 19.\par
 
% END A1.SS0.SSS0.P0.SPx3.p3
 
% BEGIN Thmtheorem40
 
\par\medskip\noindent\textbf{Lemma 40.}\par
 
% BEGIN Thmtheorem40.p1
 
Let $\sigma,\theta$ be substitutions such that $\sigma\leq_{\mathsf{ACU}}^{V}\theta$ and $\operatorname{Dom}(\sigma)\subseteq\operatorname{Dom}(\theta)$. Then $\operatorname{VRan}(\sigma)\subseteq\operatorname{Dom}(\theta)\cup%
\operatorname{VRan}(\theta)$.\par
 
% END Thmtheorem40.p1
 
% END Thmtheorem40
 
% BEGIN A1.SS0.SSS0.P0.SPx3.p4
 
Thus, given a substitution $\theta$ and an enumeration $x_{1},\ldots,x_{n}$ of $\operatorname{Dom}(\theta)\cup\operatorname{VRan}(\theta)$, we can assume (modulo $\sim_{\mathsf{ACU}}^{V}$) that any substitution in $\{\sigma\mid\sigma\leq_{\mathsf{ACU}}^{V}\theta\}$ is such that $\operatorname{Dom}(\sigma)\cup\operatorname{VRan}(\sigma)\subseteq X_{n}$. It remains to establish the contradiction mentioned above. We will do this with the help of the following result.\par
 
% END A1.SS0.SSS0.P0.SPx3.p4
 
% BEGIN Thmtheorem41
 
\par\medskip\noindent\textbf{Lemma 41.}\par
 
% BEGIN Thmtheorem41.p1
 
Let $\theta$ be a substitution, $x_{1},\ldots,x_{n}$ an enumeration of $\operatorname{Dom}(\theta)\cup\operatorname{VRan}(\theta)$, and $P\subset P_{n}$. Suppose that $\{\sigma\mid\sigma\leq_{\mathsf{ACU}}^{V}\theta\}$ contains an infinite decreasing chain $\mathfrak{s}=\sigma_{1}\sigma_{2}\ldots$ w.r.t. $>_{\mathsf{ACU}}^{V}$ such that $\sigma_{\ell+1}>_{P}\sigma_{\ell}\ \ \mbox{for all }\ell\geq 1.$ Then, there exists $(i,j)\in P_{n}\setminus P$ such that $\{\sigma\mid\sigma\leq_{\mathsf{ACU}}^{V}\theta\}$ contains an infinite decreasing chain $\tau_{1}\tau_{2}\cdots$ w.r.t. $>_{\mathsf{ACU}}^{V}$ satisfying $\tau_{p+1}>_{P\cup\{(i,j)\}}\tau_{p}\ \ \mbox{for all }p\geq 1.$\par
 
% END Thmtheorem41.p1
 
% END Thmtheorem41
 
% END A1.SS0.SSS0.P0.SPx3
 
% BEGIN A1.SS0.SSS0.P0.SPx4
 
\par\medskip\noindent\textbf{Proof.}\par
 
% BEGIN A1.SS0.SSS0.P0.SPx4.p1
 
Suppose such an infinite decreasing chain $\mathfrak{s}$ exists. We claim that there is a pair $(i,j)$ in $P_{n}\setminus P$ such that the sequence of values $c^{\sigma_{1}}_{ij},c^{\sigma_{2}}_{ij},\ldots$ is not bounded, i.e.,\par
    \[\text{for all }m\in\mathbb{N},\text{ there exists }\ell\geq 1\text{ such that:%
}\ c^{\sigma_{\ell}}_{ij}>m.\]  (3)   
If that were not the case, then there would be infinitely many substitutions in $\mathfrak{s}$ that coincide on $P_{n}\setminus P$. As a consequence, we could select substitutions $\sigma_{\ell_{1}},\sigma_{\ell_{2}}$ and $\sigma_{\ell_{3}}$ in $\mathfrak{s}$ such that:\par
  
\par\noindent  $\blacksquare$ 
% BEGIN A1.I7.i1.p1
 
$\sigma_{\ell_{3}}>_{\mathsf{ACU}}^{V}\sigma_{\ell_{2}}>_{\mathsf{ACU}}^{V}%
\sigma_{\ell_{1}}$,\par
 
% END A1.I7.i1.p1
 
 
\par\noindent  $\blacksquare$ 
% BEGIN A1.I7.i2.p1
 
$\sigma_{\ell_{1}}$ is greater than $\sigma_{\ell_{2}}$ w.r.t. $P$, and\par
 
% END A1.I7.i2.p1
 
 
\par\noindent  $\blacksquare$ 
% BEGIN A1.I7.i3.p1
 
$\sigma_{\ell_{1}}$ and $\sigma_{\ell_{3}}$ coincide on $P_{n}\setminus P$.\par
 
% END A1.I7.i3.p1
 
  
Hence, an application of Lemma 38 would yield $\sigma_{\ell_{3}}\leq_{\mathsf{ACU}}^{V}\sigma_{\ell_{2}}$, which contradicts $\sigma_{\ell_{3}}>_{\mathsf{ACU}}^{V}\sigma_{\ell_{2}}$.\par
 
% END A1.SS0.SSS0.P0.SPx4.p1
 
% BEGIN A1.SS0.SSS0.P0.SPx4.p2
 
Based on such a pair $(i,j)$, we define the infinite chain $\tau_{1}\tau_{2}\cdots$ as follows:\par
  
\par\noindent  1. 
% BEGIN A1.I8.i1.p1
 
$\tau_{1}=\sigma_{1}$.\par
 
% END A1.I8.i1.p1
 
 
\par\noindent  2. 
% BEGIN A1.I8.i2.p1
 
By (3), there is $p>1$ such that $c^{\sigma_{p}}_{ij}>c^{\sigma_{1}}_{ik}$ for all $k\in\{1,\ldots,n\}$. We choose $\tau_{2}$ as $\sigma_{p}$. The following arguments show that $\tau_{1}$ and $\tau_{2}$ satisfy the properties required of $\tau_{1}\tau_{2}\cdots$ w.r.t. $P\cup\{(i,j)\}$:\par
  
\par\noindent  $\blacksquare$ 
% BEGIN A1.I8.i2.I1.i1.p1
 
Since $\sigma_{1}>_{\mathsf{ACU}}^{V}\sigma_{p}$, we have that $\tau_{1}>_{\mathsf{ACU}}^{V}\tau_{2}$.\par
 
% END A1.I8.i2.I1.i1.p1
 
 
\par\noindent  $\blacksquare$ 
% BEGIN A1.I8.i2.I1.i2.p1
 
By selection of $p$, we know that $\sigma_{p}>_{\{(i,j)\}}\sigma_{1}$. Moreover, since $\sigma_{\ell+1}>_{P}\sigma_{\ell}$ for all $\ell\geq 1$, transitivity of $>_{P}$ yields that $\sigma_{p}>_{P}\sigma_{1}$. Consequently, we can conclude that $\tau_{2}$ is greater than $\tau_{1}$ w.r.t. $P\cup\{(i,j)\}$.\par
 
% END A1.I8.i2.I1.i2.p1
 
  
% END A1.I8.i2.p1
 
 
\par\noindent  3. 
% BEGIN A1.I8.i3.p1
 
Once we fix $\sigma_{p}$, (3) also yields $q>p$ such that $c^{\sigma_{q}}_{ij}>c^{\sigma_{p}}_{ik}$ for all $k,1\leq k\leq n$. We select $\tau_{3}$ as $\sigma_{q}$. The same arguments yield that $\tau_{2}$ and $\tau_{3}$ are as required.\par
 
% END A1.I8.i3.p1
 
 
\par\noindent  4. 
% BEGIN A1.I8.i4.p1
 
By repeating (ad infinitum) the described selection process, we can extract from $\mathfrak{s}$ an appropriate remaining sequence of substitutions $\tau_{4}\tau_{5}\cdots$.\par
 
% END A1.I8.i4.p1
 
  
The described process ensures that each selected substitution $\tau_{\ell}$ ($\ell\geq 1$) belongs to $\{\sigma\mid\sigma\leq_{\mathsf{ACU}}^{V}\theta\}$. Thus, $\tau_{1}\tau_{2}\cdots$ is an infinite decreasing chain satisfying the claim of the lemma. $\hfill\blacktriangleleft$\par
 
% END A1.SS0.SSS0.P0.SPx4.p2
 
% BEGIN A1.SS0.SSS0.P0.SPx4.p3
 
Finally, by using the previous lemma, we can show the main result of this section.\par
 
% END A1.SS0.SSS0.P0.SPx4.p3
 
% BEGIN Thmtheorem42
 
\par\medskip\noindent\textbf{Lemma 42.}\par
 
% BEGIN Thmtheorem42.p1
 
For a given substitution $\theta$, the set $\{\sigma\mid\sigma\leq_{\mathsf{ACU}}^{V}\theta\}$ does not contain an infinite decreasing chain w.r.t. $>_{\mathsf{ACU}}^{V}$.\par
 
% END Thmtheorem42.p1
 
% END Thmtheorem42
 
% END A1.SS0.SSS0.P0.SPx4
 
% BEGIN A1.SS0.SSS0.P0.SPx5
 
\par\medskip\noindent\textbf{Proof.}\par
 
% BEGIN A1.SS0.SSS0.P0.SPx5.p1
 
Suppose $\{\sigma\mid\sigma\leq_{\mathsf{ACU}}^{V}\theta\}$ contains an infinite decreasing chain $\mathfrak{s}$ w.r.t. $>_{\mathsf{ACU}}^{V}$. By Definition 36, we have $\eta>_{\emptyset}\eta^{\prime}$ for all $\eta,\eta^{\prime}\in\{\sigma\mid\sigma\leq_{\mathsf{ACU}}^{V}\theta\}$. This means that $\mathfrak{s}$ and $P=\emptyset$ satisfy the hypothesis of Lemma 41. Hence, there is $(i,j)\in P_{n}$ such that $\{\sigma\mid\sigma\leq_{\mathsf{ACU}}^{V}\theta\}$ contains an infinite decreasing chain $\tau_{1}\tau_{2}\cdots$ satisfying $\tau_{p+1}>_{\{(i,j)\}}\tau_{p}$ for all $p\geq 1$. This new chain now satisfies the hypothesis of Lemma 41 w.r.t. $P=\{(i,j)\}$. Thus, by a sequence of $n^{2}-1$ further applications of Lemma 41, we conclude that $\{\sigma\mid\sigma\leq_{\mathsf{ACU}}^{V}\theta\}$ contains an infinite decreasing chain $\sigma_{1}\sigma_{2}\cdots$ such that $\sigma_{p+1}>_{P_{n}}\sigma_{p}$ for all $p\geq 1$. Finally, it is not hard to verify that this last chain must contain three substitutions $\sigma_{\ell_{1}},\sigma_{\ell_{2}}$ and $\sigma_{\ell_{3}}$ such that\par
  
\par\noindent  $\blacksquare$ 
% BEGIN A1.I9.i1.p1
 
$\sigma_{\ell_{3}}>_{\mathsf{ACU}}^{V}\sigma_{\ell_{2}}>_{\mathsf{ACU}}^{V}%
\sigma_{\ell_{1}}$,\par
 
% END A1.I9.i1.p1
 
 
\par\noindent  $\blacksquare$ 
% BEGIN A1.I9.i2.p1
 
$\sigma_{\ell_{1}}$ is greater than $\sigma_{\ell_{2}}$ w.r.t. $P=P_{n}$, and\par
 
% END A1.I9.i2.p1
 
 
\par\noindent  $\blacksquare$ 
% BEGIN A1.I9.i3.p1
 
$\sigma_{\ell_{1}}$ and $\sigma_{\ell_{3}}$ coincide on $P_{n}\setminus P=\emptyset$.\par
 
% END A1.I9.i3.p1
 
  
Hence, the application of Lemma 38 yields that $\sigma_{\ell_{3}}\leq_{\mathsf{ACU}}^{V}\sigma_{\ell_{2}}$, which contradicts the existence of $\sigma_{1}\sigma_{2}\cdots$ since $\sigma_{\ell_{3}}>_{\mathsf{ACU}}^{V}\sigma_{\ell_{2}}$. Thus, we have derived a contradiction from our initial assumption. This concludes the proof of the lemma. $\hfill\blacktriangleleft$\par
 
% END A1.SS0.SSS0.P0.SPx5.p1
 
% BEGIN A1.SS0.SSS0.P0.SPx5.p2
 
Lemma 42 shows that the unrestricted unification type of $\mathsf{ACU}$ cannot be zero. In fact, this lemma implies that every unifier $\theta$ is above a minimal unifier, and thus the set of minimal unifiers is complete. By Theorem 3, we can conclude that the unification type cannot be zero, which completes the proof of Theorem 31.\par
 
% END A1.SS0.SSS0.P0.SPx5.p2
 
% BEGIN A1.SS0.SSS0.P0.SPx5.p3
 
\par
 
% END A1.SS0.SSS0.P0.SPx5.p3
 
% END A1.SS0.SSS0.P0.SPx5
 
% END A1
